% concatenated from l4dc2026-sample.tex, math_commands.tex
\documentclass{l4dc2026}
% If your paper is accepted, please use the "final" option
% for the camera-ready version
% \documentclass[final]{l4dc2026}

% The following packages will be automatically loaded:
% amsmath, amssymb, natbib, graphicx, url, algorithm2e

% \title[Short Title]{Safety Beyond the Training Data: Robust Planning with Learned Dynamics via Conformal Prediction and System Level Synthesis}
% \title[Short Title]{Robust Out-of-Distribution Planning and Control with Learned Dynamics via Conformal Prediction and System Level Synthesis}
\title[Robust Out-of-Distribution MPC via Conformalized System Level Synthesis]{Safety Beyond the Training Data: Robust Out-of-Distribution MPC via Conformalized System Level Synthesis}
\usepackage{times}
\usepackage{graphics}
\usepackage{sidecap}
\usepackage{booktabs}
\usepackage{multirow}
%%%%% NEW MATH DEFINITIONS %%%%%

\usepackage{amsmath,amsfonts,bm}

% Mark sections of captions for referring to divisions of figures
\newcommand{\figleft}{{\em (Left)}}
\newcommand{\figcenter}{{\em (Center)}}
\newcommand{\figright}{{\em (Right)}}
\newcommand{\figtop}{{\em (Top)}}
\newcommand{\figbottom}{{\em (Bottom)}}
\newcommand{\captiona}{{\em (a)}}
\newcommand{\captionb}{{\em (b)}}
\newcommand{\captionc}{{\em (c)}}
\newcommand{\captiond}{{\em (d)}}

% Highlight a newly defined term
\newcommand{\newterm}[1]{{\bf #1}}


% Figure reference, lower-case.
\def\figref#1{figure~\ref{#1}}
% Figure reference, capital. For start of sentence
\def\Figref#1{Figure~\ref{#1}}
\def\twofigref#1#2{figures \ref{#1} and \ref{#2}}
\def\quadfigref#1#2#3#4{figures \ref{#1}, \ref{#2}, \ref{#3} and \ref{#4}}
% Section reference, lower-case.
\def\secref#1{section~\ref{#1}}
% Section reference, capital.
\def\Secref#1{Section~\ref{#1}}
% Reference to two sections.
\def\twosecrefs#1#2{sections \ref{#1} and \ref{#2}}
% Reference to three sections.
\def\secrefs#1#2#3{sections \ref{#1}, \ref{#2} and \ref{#3}}
% Reference to an equation, lower-case.
% \def\eqref#1{equation~\ref{#1}}
% Reference to an equation, upper case
\def\Eqref#1{Equation~\ref{#1}}
% A raw reference to an equation---avoid using if possible
\def\plaineqref#1{\ref{#1}}
% Reference to a chapter, lower-case.
\def\chapref#1{chapter~\ref{#1}}
% Reference to an equation, upper case.
\def\Chapref#1{Chapter~\ref{#1}}
% Reference to a range of chapters
\def\rangechapref#1#2{chapters\ref{#1}--\ref{#2}}
% Reference to an algorithm, lower-case.
\def\algref#1{algorithm~\ref{#1}}
% Reference to an algorithm, upper case.
\def\Algref#1{Algorithm~\ref{#1}}
\def\twoalgref#1#2{algorithms \ref{#1} and \ref{#2}}
\def\Twoalgref#1#2{Algorithms \ref{#1} and \ref{#2}}
% Reference to a part, lower case
\def\partref#1{part~\ref{#1}}
% Reference to a part, upper case
\def\Partref#1{Part~\ref{#1}}
\def\twopartref#1#2{parts \ref{#1} and \ref{#2}}

% \def\ceil#1{\lceil #1 \rceil}
% \def\floor#1{\lfloor #1 \rfloor}
\def\1{\bm{1}}
\newcommand{\train}{\mathcal{D}_{\mathrm{train}}}
\newcommand{\valid}{\mathcal{D}_{\mathrm{valid}}}
\newcommand{\calib}{\mathcal{D}_{\mathrm{calib}}}
\newcommand{\test}{\mathcal{D}_{\mathrm{test}}}
\newcommand{\tvd}{d_\textrm{TV}}

\def\eps{{\epsilon}}


% Random variables
\def\reta{{\textnormal{$\eta$}}}
\def\ra{{\textnormal{a}}}
\def\rb{{\textnormal{b}}}
\def\rc{{\textnormal{c}}}
\def\rd{{\textnormal{d}}}
\def\re{{\textnormal{e}}}
\def\rf{{\textnormal{f}}}
\def\rg{{\textnormal{g}}}
\def\rh{{\textnormal{h}}}
\def\ri{{\textnormal{i}}}
\def\rj{{\textnormal{j}}}
\def\rk{{\textnormal{k}}}
\def\rl{{\textnormal{l}}}
% rm is already a command, just don't name any random variables m
\def\rn{{\textnormal{n}}}
\def\ro{{\textnormal{o}}}
\def\rp{{\textnormal{p}}}
\def\rq{{\textnormal{q}}}
\def\rr{{\textnormal{r}}}
\def\rs{{\textnormal{s}}}
\def\rt{{\textnormal{t}}}
\def\ru{{\textnormal{u}}}
\def\rv{{\textnormal{v}}}
\def\rw{{\textnormal{w}}}
\def\rx{{\textnormal{x}}}
\def\ry{{\textnormal{y}}}
\def\rz{{\textnormal{z}}}

% Random vectors
\def\rvepsilon{{\mathbf{\epsilon}}}
\def\rvtheta{{\mathbf{\theta}}}
\def\rva{{\mathbf{a}}}
\def\rvb{{\mathbf{b}}}
\def\rvc{{\mathbf{c}}}
\def\rvd{{\mathbf{d}}}
\def\rve{{\mathbf{e}}}
\def\rvf{{\mathbf{f}}}
\def\rvg{{\mathbf{g}}}
\def\rvh{{\mathbf{h}}}
\def\rvu{{\mathbf{i}}}
\def\rvj{{\mathbf{j}}}
\def\rvk{{\mathbf{k}}}
\def\rvl{{\mathbf{l}}}
\def\rvm{{\mathbf{m}}}
\def\rvn{{\mathbf{n}}}
\def\rvo{{\mathbf{o}}}
\def\rvp{{\mathbf{p}}}
\def\rvq{{\mathbf{q}}}
\def\rvr{{\mathbf{r}}}
\def\rvs{{\mathbf{s}}}
\def\rvt{{\mathbf{t}}}
\def\rvu{{\mathbf{u}}}
\def\rvv{{\mathbf{v}}}
\def\rvw{{\mathbf{w}}}
\def\rvx{{\mathbf{x}}}
\def\rvy{{\mathbf{y}}}
\def\rvz{{\mathbf{z}}}

% Elements of random vectors
\def\erva{{\textnormal{a}}}
\def\ervb{{\textnormal{b}}}
\def\ervc{{\textnormal{c}}}
\def\ervd{{\textnormal{d}}}
\def\erve{{\textnormal{e}}}
\def\ervf{{\textnormal{f}}}
\def\ervg{{\textnormal{g}}}
\def\ervh{{\textnormal{h}}}
\def\ervi{{\textnormal{i}}}
\def\ervj{{\textnormal{j}}}
\def\ervk{{\textnormal{k}}}
\def\ervl{{\textnormal{l}}}
\def\ervm{{\textnormal{m}}}
\def\ervn{{\textnormal{n}}}
\def\ervo{{\textnormal{o}}}
\def\ervp{{\textnormal{p}}}
\def\ervq{{\textnormal{q}}}
\def\ervr{{\textnormal{r}}}
\def\ervs{{\textnormal{s}}}
\def\ervt{{\textnormal{t}}}
\def\ervu{{\textnormal{u}}}
\def\ervv{{\textnormal{v}}}
\def\ervw{{\textnormal{w}}}
\def\ervx{{\textnormal{x}}}
\def\ervy{{\textnormal{y}}}
\def\ervz{{\textnormal{z}}}

% Random matrices
\def\rmA{{\mathbf{A}}}
\def\rmB{{\mathbf{B}}}
\def\rmC{{\mathbf{C}}}
\def\rmD{{\mathbf{D}}}
\def\rmE{{\mathbf{E}}}
\def\rmF{{\mathbf{F}}}
\def\rmG{{\mathbf{G}}}
\def\rmH{{\mathbf{H}}}
\def\rmI{{\mathbf{I}}}
\def\rmJ{{\mathbf{J}}}
\def\rmK{{\mathbf{K}}}
\def\rmL{{\mathbf{L}}}
\def\rmM{{\mathbf{M}}}
\def\rmN{{\mathbf{N}}}
\def\rmO{{\mathbf{O}}}
\def\rmP{{\mathbf{P}}}
\def\rmQ{{\mathbf{Q}}}
\def\rmR{{\mathbf{R}}}
\def\rmS{{\mathbf{S}}}
\def\rmT{{\mathbf{T}}}
\def\rmU{{\mathbf{U}}}
\def\rmV{{\mathbf{V}}}
\def\rmW{{\mathbf{W}}}
\def\rmX{{\mathbf{X}}}
\def\rmY{{\mathbf{Y}}}
\def\rmZ{{\mathbf{Z}}}

% Elements of random matrices
\def\ermA{{\textnormal{A}}}
\def\ermB{{\textnormal{B}}}
\def\ermC{{\textnormal{C}}}
\def\ermD{{\textnormal{D}}}
\def\ermE{{\textnormal{E}}}
\def\ermF{{\textnormal{F}}}
\def\ermG{{\textnormal{G}}}
\def\ermH{{\textnormal{H}}}
\def\ermI{{\textnormal{I}}}
\def\ermJ{{\textnormal{J}}}
\def\ermK{{\textnormal{K}}}
\def\ermL{{\textnormal{L}}}
\def\ermM{{\textnormal{M}}}
\def\ermN{{\textnormal{N}}}
\def\ermO{{\textnormal{O}}}
\def\ermP{{\textnormal{P}}}
\def\ermQ{{\textnormal{Q}}}
\def\ermR{{\textnormal{R}}}
\def\ermS{{\textnormal{S}}}
\def\ermT{{\textnormal{T}}}
\def\ermU{{\textnormal{U}}}
\def\ermV{{\textnormal{V}}}
\def\ermW{{\textnormal{W}}}
\def\ermX{{\textnormal{X}}}
\def\ermY{{\textnormal{Y}}}
\def\ermZ{{\textnormal{Z}}}

% Vectors
\def\vzero{{\bm{0}}}
\def\vone{{\bm{1}}}
\def\vmu{{\bm{\mu}}}
\def\vtheta{{\bm{\theta}}}
\def\va{{\bm{a}}}
\def\vb{{\bm{b}}}
\def\vc{{\bm{c}}}
\def\vd{{\bm{d}}}
\def\ve{{\bm{e}}}
\def\vf{{\bm{f}}}
\def\vg{{\bm{g}}}
\def\vh{{\bm{h}}}
\def\vi{{\bm{i}}}
\def\vj{{\bm{j}}}
\def\vk{{\bm{k}}}
\def\vl{{\bm{l}}}
\def\vm{{\bm{m}}}
\def\vn{{\bm{n}}}
\def\vo{{\bm{o}}}
\def\vp{{\bm{p}}}
\def\vq{{\bm{q}}}
\def\vr{{\bm{r}}}
\def\vs{{\bm{s}}}
\def\vt{{\bm{t}}}
\def\vu{{\bm{u}}}
\def\vv{{\bm{v}}}
\def\vw{{\bm{w}}}
\def\vx{{\bm{x}}}
\def\vy{{\bm{y}}}
\def\vz{{\bm{z}}}

% Elements of vectors
\def\evalpha{{\alpha}}
\def\evbeta{{\beta}}
\def\evepsilon{{\epsilon}}
\def\evlambda{{\lambda}}
\def\evomega{{\omega}}
\def\evmu{{\mu}}
\def\evpsi{{\psi}}
\def\evsigma{{\sigma}}
\def\evtheta{{\theta}}
\def\eva{{a}}
\def\evb{{b}}
\def\evc{{c}}
\def\evd{{d}}
\def\eve{{e}}
\def\evf{{f}}
\def\evg{{g}}
\def\evh{{h}}
\def\evi{{i}}
\def\evj{{j}}
\def\evk{{k}}
\def\evl{{l}}
\def\evm{{m}}
\def\evn{{n}}
\def\evo{{o}}
\def\evp{{p}}
\def\evq{{q}}
\def\evr{{r}}
\def\evs{{s}}
\def\evt{{t}}
\def\evu{{u}}
\def\evv{{v}}
\def\evw{{w}}
\def\evx{{x}}
\def\evy{{y}}
\def\evz{{z}}

% Matrix
\def\mA{{\bm{A}}}
\def\mB{{\bm{B}}}
\def\mC{{\bm{C}}}
\def\mD{{\bm{D}}}
\def\mE{{\bm{E}}}
\def\mF{{\bm{F}}}
\def\mG{{\bm{G}}}
\def\mH{{\bm{H}}}
\def\mI{{\bm{I}}}
\def\mJ{{\bm{J}}}
\def\mK{{\bm{K}}}
\def\mL{{\bm{L}}}
\def\mM{{\bm{M}}}
\def\mN{{\bm{N}}}
\def\mO{{\bm{O}}}
\def\mP{{\bm{P}}}
\def\mQ{{\bm{Q}}}
\def\mR{{\bm{R}}}
\def\mS{{\bm{S}}}
\def\mT{{\bm{T}}}
\def\mU{{\bm{U}}}
\def\mV{{\bm{V}}}
\def\mW{{\bm{W}}}
\def\mX{{\bm{X}}}
\def\mY{{\bm{Y}}}
\def\mZ{{\bm{Z}}}
\def\mBeta{{\bm{\beta}}}
\def\mPhi{{\bm{\Phi}}}
\def\mLambda{{\bm{\Lambda}}}
\def\mSigma{{\bm{\Sigma}}}

% Tensor
\DeclareMathAlphabet{\mathsfit}{\encodingdefault}{\sfdefault}{m}{sl}
\SetMathAlphabet{\mathsfit}{bold}{\encodingdefault}{\sfdefault}{bx}{n}
\newcommand{\tens}[1]{\bm{\mathsfit{#1}}}
\def\tA{{\tens{A}}}
\def\tB{{\tens{B}}}
\def\tC{{\tens{C}}}
\def\tD{{\tens{D}}}
\def\tE{{\tens{E}}}
\def\tF{{\tens{F}}}
\def\tG{{\tens{G}}}
\def\tH{{\tens{H}}}
\def\tI{{\tens{I}}}
\def\tJ{{\tens{J}}}
\def\tK{{\tens{K}}}
\def\tL{{\tens{L}}}
\def\tM{{\tens{M}}}
\def\tN{{\tens{N}}}
\def\tO{{\tens{O}}}
\def\tP{{\tens{P}}}
\def\tQ{{\tens{Q}}}
\def\tR{{\tens{R}}}
\def\tS{{\tens{S}}}
\def\tT{{\tens{T}}}
\def\tU{{\tens{U}}}
\def\tV{{\tens{V}}}
\def\tW{{\tens{W}}}
\def\tX{{\tens{X}}}
\def\tY{{\tens{Y}}}
\def\tZ{{\tens{Z}}}


% Graph
\def\gA{{\mathcal{A}}}
\def\gB{{\mathcal{B}}}
\def\gC{{\mathcal{C}}}
\def\gD{{\mathcal{D}}}
\def\gE{{\mathcal{E}}}
\def\gF{{\mathcal{F}}}
\def\gG{{\mathcal{G}}}
\def\gH{{\mathcal{H}}}
\def\gI{{\mathcal{I}}}
\def\gJ{{\mathcal{J}}}
\def\gK{{\mathcal{K}}}
\def\gL{{\mathcal{L}}}
\def\gM{{\mathcal{M}}}
\def\gN{{\mathcal{N}}}
\def\gO{{\mathcal{O}}}
\def\gP{{\mathcal{P}}}
\def\gQ{{\mathcal{Q}}}
\def\gR{{\mathcal{R}}}
\def\gS{{\mathcal{S}}}
\def\gT{{\mathcal{T}}}
\def\gU{{\mathcal{U}}}
\def\gV{{\mathcal{V}}}
\def\gW{{\mathcal{W}}}
\def\gX{{\mathcal{X}}}
\def\gY{{\mathcal{Y}}}
\def\gZ{{\mathcal{Z}}}

% Sets
\def\sA{{\mathbb{A}}}
\def\sB{{\mathbb{B}}}
\def\sC{{\mathbb{C}}}
\def\sD{{\mathbb{D}}}
% Don't use a set called E, because this would be the same as our symbol
% for expectation.
\def\sF{{\mathbb{F}}}
\def\sG{{\mathbb{G}}}
\def\sH{{\mathbb{H}}}
\def\sI{{\mathbb{I}}}
\def\sJ{{\mathbb{J}}}
\def\sK{{\mathbb{K}}}
\def\sL{{\mathbb{L}}}
\def\sM{{\mathbb{M}}}
\def\sN{{\mathbb{N}}}
\def\sO{{\mathbb{O}}}
\def\sP{{\mathbb{P}}}
\def\sQ{{\mathbb{Q}}}
\def\sR{{\mathbb{R}}}
\def\sS{{\mathbb{S}}}
\def\sT{{\mathbb{T}}}
\def\sU{{\mathbb{U}}}
\def\sV{{\mathbb{V}}}
\def\sW{{\mathbb{W}}}
\def\sX{{\mathbb{X}}}
\def\sY{{\mathbb{Y}}}
\def\sZ{{\mathbb{Z}}}

% Entries of a matrix
\def\emLambda{{\Lambda}}
\def\emA{{A}}
\def\emB{{B}}
\def\emC{{C}}
\def\emD{{D}}
\def\emE{{E}}
\def\emF{{F}}
\def\emG{{G}}
\def\emH{{H}}
\def\emI{{I}}
\def\emJ{{J}}
\def\emK{{K}}
\def\emL{{L}}
\def\emM{{M}}
\def\emN{{N}}
\def\emO{{O}}
\def\emP{{P}}
\def\emQ{{Q}}
\def\emR{{R}}
\def\emS{{S}}
\def\emT{{T}}
\def\emU{{U}}
\def\emV{{V}}
\def\emW{{W}}
\def\emX{{X}}
\def\emY{{Y}}
\def\emZ{{Z}}
\def\emSigma{{\Sigma}}

% entries of a tensor
% Same font as tensor, without \bm wrapper
\newcommand{\etens}[1]{\mathsfit{#1}}
\def\etLambda{{\etens{\Lambda}}}
\def\etA{{\etens{A}}}
\def\etB{{\etens{B}}}
\def\etC{{\etens{C}}}
\def\etD{{\etens{D}}}
\def\etE{{\etens{E}}}
\def\etF{{\etens{F}}}
\def\etG{{\etens{G}}}
\def\etH{{\etens{H}}}
\def\etI{{\etens{I}}}
\def\etJ{{\etens{J}}}
\def\etK{{\etens{K}}}
\def\etL{{\etens{L}}}
\def\etM{{\etens{M}}}
\def\etN{{\etens{N}}}
\def\etO{{\etens{O}}}
\def\etP{{\etens{P}}}
\def\etQ{{\etens{Q}}}
\def\etR{{\etens{R}}}
\def\etS{{\etens{S}}}
\def\etT{{\etens{T}}}
\def\etU{{\etens{U}}}
\def\etV{{\etens{V}}}
\def\etW{{\etens{W}}}
\def\etX{{\etens{X}}}
\def\etY{{\etens{Y}}}
\def\etZ{{\etens{Z}}}

% The true underlying data generating distribution
\newcommand{\pdata}{p_{\rm{data}}}
% The empirical distribution defined by the training set
\newcommand{\ptrain}{\hat{p}_{\rm{data}}}
\newcommand{\Ptrain}{\hat{P}_{\rm{data}}}
% The model distribution
\newcommand{\pmodel}{p_{\rm{model}}}
\newcommand{\Pmodel}{P_{\rm{model}}}
\newcommand{\ptildemodel}{\tilde{p}_{\rm{model}}}
% Stochastic autoencoder distributions
\newcommand{\pencode}{p_{\rm{encoder}}}
\newcommand{\pdecode}{p_{\rm{decoder}}}
\newcommand{\precons}{p_{\rm{reconstruct}}}

\newcommand{\laplace}{\mathrm{Laplace}} % Laplace distribution

\newcommand{\E}{\mathbb{E}}
\newcommand{\Ls}{\mathcal{L}}
\newcommand{\R}{\mathbb{R}}
\newcommand{\emp}{\tilde{p}}
\newcommand{\lr}{\alpha}
\newcommand{\reg}{\lambda}
\newcommand{\rect}{\mathrm{rectifier}}
\newcommand{\softmax}{\mathrm{softmax}}
\newcommand{\sigmoid}{\sigma}
\newcommand{\softplus}{\zeta}
\newcommand{\KL}{D_{\mathrm{KL}}}
\newcommand{\Var}{\mathrm{Var}}
\newcommand{\standarderror}{\mathrm{SE}}
\newcommand{\Cov}{\mathrm{Cov}}
% Wolfram Mathworld says $L^2$ is for function spaces and $\ell^2$ is for vectors
% But then they seem to use $L^2$ for vectors throughout the site, and so does
% wikipedia.
\newcommand{\normlzero}{L^0}
\newcommand{\normlone}{L^1}
\newcommand{\normltwo}{L^2}
\newcommand{\normlp}{L^p}
\newcommand{\normmax}{L^\infty}

\newcommand{\parents}{Pa} % See usage in notation.tex. Chosen to match Daphne's book.

\DeclareMathOperator*{\argmax}{arg\,max}
\DeclareMathOperator*{\argmin}{arg\,min}

\DeclareMathOperator{\sign}{sign}
\DeclareMathOperator{\Tr}{Tr}
\let\ab\allowbreak

\usepackage{multicol}
\usepackage{wrapfig}
\usepackage{algorithm,algpseudocode}
% \setcitestyle{numbers,sort&compress,square}
\usepackage{caption}
\usepackage{subcaption}

\newcommand{\glen}[1]{{\textcolor[rgb]{1,0,0}{\textsf{[Glen: {#1}]}}}}
\newcommand{\antoine}[1]{{\textcolor[rgb]{1,0,0}{\textsf{[Antoine: {#1}]}}}}

\newcommand{\defeq}{:=}
\newcommand{\eqdef}{=:}
\newcommand{\anutamchange}[1]{{\color{orange}#1}}
% remove vspace temporarily
%\renewcommand{\vspace}[1]{}

\newenvironment{proofsketch}{%
  \renewcommand{\proofname}{Proof Sketch}\proof}{\endproof}
\let\lamon\l
\renewcommand{\l}{_\textrm{lin}}


\newcommand{\methodnospace}{CP-SLS-MPC}
\newcommand{\method}{CP-SLS-MPC }
\newcommand{\xz}{\bar x}
\newcommand{\Z}{\mathbf{z}}
\newcommand{\V}{\mathbf{v}}
\newcommand{\U}{\mathbf{u}}
\newcommand{\M}{M}


\renewcommand{\P}{\mathbf \Phi}
\newcommand{\Y}{\mathbf{y}}
\newcommand{\X}{\mathbf{x}}

\newcommand{\nx}{n_\mathrm{x}}
\renewcommand{\nu}{n_\mathrm{u}}
\newcommand{\nc}{n_\mathrm{c}}
\newcommand{\nf}{n_\mathrm{f}}


\newcommand{\nw}{n_\mathrm{w}}
\newcommand{\ny}{{n_\mathrm{y}}}

\newcommand{\Qf}{P}

\newcommand{\W}{\mathbf{w}}
\newcommand{\f}{\mathrm{f}}
\newcommand{\covariance}{\mathbf{\Sigma}}

\newcommand{\T}{^\top}

\newcommand{\x}{\mathrm{x}}
\renewcommand{\u}{\mathrm{u}}
\newcommand{\w}{\mathrm{w}}

\newcommand{\B}{ \mathcal{B}_\infty}
\newcommand{\mathdef}{=\vcentcolon}
\newcommand{\red}[1]{{\color{red}{#1}}}


\newcommand{\Px}{\bm{\Phi}^{\mathrm{x}}}
\newcommand{\Pu}{\bm{\Phi}^{\mathrm{u}}}
\newcommand{\F}{\mathcal{F}}




% Use \Name{Author Name} to specify the name.
% If the surname contains spaces, enclose the surname
% in braces, e.g. \Name{John {Smith Jones}} similarly
% if the name has a "von" part, e.g \Name{Jane {de Winter}}.
% If the first letter in the forenames is a diacritic
% enclose the diacritic in braces, e.g. \Name{{\'E}louise Smith}

% Two authors with the same address
% \coltauthor{\Name{Author Name1} \Email{abc@sample.com}\and
%  \Name{Author Name2} \Email{xyz@sample.com}\\
%  \addr Address}

% Three or more authors with the same address:
% \coltauthor{\Name{Author Name1} \Email{an1@sample.com}\\
%  \Name{Author Name2} \Email{an2@sample.com}\\
%  \Name{Author Name3} \Email{an3@sample.com}\\
%  \addr Address}

% Authors with different addresses:
% \author{%
%  \Name{Anutam Srinivasan}$^1$ \Email{asrinivasan350@gatech.edu}\\
%  \Name{Glen Chou}$^2$ \Email{chou@gatech.edu}\\
%  \addr Georgia Institute of Technology. $^1$ School of ECE, $^2$ Schools of Cybersecurity \& Privacy \& Aerospace Eng.%
% }
% \author{\Name{Anutam Srinivasan}$~^1$ \and
%  \Name{Glen Chou}$~^2$ \Email{\{asrinivasan350, chou\}@gatech.edu}\\
%  \addr Georgia Institute of Technology. $^1$ School of ECE, $^2$ Schools of Cybersecurity \& Privacy \& Aerospace Eng.%}
%  }
\author{%
 \Name{Anutam Srinivasan}$^{1}$ \Email{asrinivasan350@gatech.edu}\\
 \Name{Antoine Leeman}$^{2}$ \Email{aleeman@ethz.ch}\\
 \Name{Glen Chou}$^1$ \Email{chou@gatech.edu}\\
 \addr $^1$Georgia Institute of Technology. $^2$ ETH Zurich.%
}
\usepackage{xcolor}
\renewcommand{\red}[1]{{\color{red}{#1}}}


\begin{document}

\maketitle
\vspace{-20pt}
\begin{abstract}%
We present a novel framework for robust out-of-distribution planning and control using conformal prediction (CP) and system level synthesis (SLS), addressing the challenge of ensuring safety and robustness when using learned dynamics models beyond the training data distribution. We first derive high-confidence model error bounds using weighted CP with a learned, state-control-dependent covariance model. These bounds are integrated into an SLS-based robust nonlinear model predictive control (MPC) formulation, which performs constraint tightening over the prediction horizon via volume-optimized forward reachable sets. We provide theoretical guarantees on coverage and robustness under distributional drift, and analyze the impact of data density and trajectory tube size on prediction coverage. Empirically, we demonstrate our method on nonlinear systems of increasing complexity, including a 4D car and a {12D} quadcopter, improving safety and robustness compared to fixed-bound and non-robust baselines, especially outside of the data distribution.
\end{abstract}

\begin{keywords}%
  safe learning-based control, conformal prediction, MPC, system level synthesis%
\end{keywords}

\vspace{-12pt}
\section{Introduction}
\vspace{-3pt}

% % Provably-safe motion planning for systems with unknown dynamics is essential for real-world robot deployment. While traditional planning algorithms perform reliably when the system dynamics are precisely known, the dynamics are often uncertain or inaccurately modeled in practice for many robotic systems. Data-driven approaches, such as model-based reinforcement learning \cite{DBLP:journals/ftml/MoerlandBPJ23}, attempt to learn the dynamics from data and plan accordingly. However, these methods can be unreliable and unsafe because planners may exploit inaccuracies in the learned model, producing trajectories that the real system cannot follow, which can result in unpredictable or hazardous behavior \cite{DBLP:conf/cdc/ChouOB21}. This motivates uncertainty-constrained planners that can ensure safety by quantifying the potential tracking error of the true system relative to the planned trajectory and using this bound to design trajectories that can be robustly executed.
% Provably-safe planning under unknown dynamics is crucial for real-world robotics. This poses difficulties for classical planners, which assume precise models; in practice, model error and uncertainty are unavoidable.
% %Data-driven methods such as model-based reinforcement learning (RL) can learn dynamics from data but may exploit model inaccuracies, leading to unsafe or infeasible trajectories. 
% This motivate uncertainty-constrained planners that quantify tracking error between true and learned dynamics to ensure robust, safe execution.
% % Provably safe motion planning algorithms for unknown systems are critical for deploying robots in the real world. While planners are reliable when the system dynamics are known exactly, the dynamics often may be poorly modeled or unknown. To address this, data-driven methods (i.e. model-based reinforcement learning) learn the dynamics from data and plan with the learned model. However, such methods can be unsafe, in part because the planner can and will exploit errors in the learned dynamics to return trajectories which cannot actually be tracked on the real system, leading to unpredictable, unsafe behavior when executed. Thus, to guarantee safety, it is of major interest to establish a bound on the error that the true system may see when attempting to track a trajectory planned with the learned dynamics, and to use it to guide the planning of robustly-trackable trajectories.
% % In recent years, several methods have been proposed to enable safe planning with learned dynamics \cite{DBLP:conf/cdc/ChouOB21, DBLP:journals/ral/KnuthCOB21, DBLP:conf/icra/KnuthCRM23, wasiela2024learning}. These approaches typically estimate a model error bound valid over a subset of the state-control space and constrain the closed-loop behavior to remain within this ``trusted domain." They leverage closed-loop reachability analysis based on contraction theory to efficiently compute reachable tubes that account for spatially varying model error and tracking of distinct reference trajectories. 
% % Recent methods have enabled safe planning with learned dynamics \cite{DBLP:journals/ral/KnuthCOB21, DBLP:conf/icra/KnuthCRM23, wasiela2024learning} by estimating model error bounds over subsets of the state-control space and constraining closed-loop behavior within this ``trusted domain.” They use closed-loop reachability analysis to compute reachable tubes that account for spatially varying model error and tracking of reference trajectories. 
% % However, the trusted domain is usually defined as a bounded neighborhood around the training data, which limits the use of learned dynamics far from the data and prevents planning when start and goal states lie in disconnected regions of the ``trusted domain". This is challenging in practical robotics, where data is scarce and exhaustive coverage is infeasible. 
% % Reaching a target may therefore require temporarily or entirely operating out of distribution (OOD), tolerating degraded performance while maintaining safety.
% Recent methods enable safe planning with learned dynamics \citep{DBLP:journals/ral/KnuthCOB21, DBLP:conf/icra/KnuthCRM23, wasiela2024learning} by estimating spatially varying model error bounds and constraining closed-loop behavior within a ``trusted domain" via reachable tubes.
% However, the trusted domain is typically a bounded neighborhood around the training data, which prevents plans that stray far from the data. 
% In practice, data scarcity in many tasks often forces robots to operate temporarily out of distribution (OOD), tolerating degraded performance while maintaining safety.
% Moreover, these methods are inefficient in their search procedure, and do not involve a search process that actively reduces model uncertainty.
% %A key gap is a fast approach that actively shrinks uncertainty to enable less conservative, uncertainty-informed planning.
% % These approaches are also inefficient, relying on sampling-based planners that do not actively reduce model uncertainty to safely reach the goal. A key missing component is an optimization-based formulation and an error metric to shrink uncertainty and guide trajectory design. Moreover, fast optimization-driven methods would allow real-time replanning and online adaptation of model error bounds, enabling less conservative and more informed planning.
% % Moreover, these methods are computationally slow \cite{DBLP:conf/cdc/ChouOB21}, relying on kinodynamic sampling-based planners that perform wasteful reachable tube–environment collision checks while exploring toward the goal in an uninformed manner. A key component missing in these approaches is an efficient optimization-based formulation and an error metric that can actively guide trajectory design. Faster, optimization-driven methods would enable real-time replanning and online adaptation of model error bounds as new data becomes available, resulting in less conservative, more informed planning.
% % In recent years, several methods have been proposed to enable safe planning with learned dynamics \cite{}. These approaches generally estimate a model error bound that is valid over a subset of the state-control space and constrain the closed-loop behavior to remain within this ``trusted set". These methods leverage closed-loop reachability analysis based on contraction theory to efficiently compute reachable sets that account for spatially-varying model error and when tracking distinct reference trajectories. However, these methods define the trusted set to be a bounded neighborhoods around the training data, which inherently prevents use of the learned dynamics sufficiently far from data and prevents planning when the start and goal states lie in disconnected regions of the validated set. This poses challenges for practical robotics applications, where data is often scarce, and exhaustive coverage of the state-control space is infeasible. Reaching a target may therefore require temporarily operating out of distribution—or even entirely out of distribution—while tolerating degraded performance, provided safety can still be ensured.
% % In recent years, methods have been proposed to plan safely with learned dynamics. Broadly speaking, they estimate a model error bound valid over a subset of state-control space and restrict the closed-loop behavior to remain in this set. Closed-loop reachability analysis is done via contraction theory, which can efficiently compute reachable sets for tracking different reference trajectories. However, these methods are typically limited to a neighborhood around the training data, which prevents planning when the start and goal states lie in disconnected regions. In practical robotics applications, data is often limited, and exhaustive coverage of the state-control space is infeasible. Reaching a target may require temporarily operating out of distribution, or even operating entirely out of distribution while accepting degraded performance, as long as safety can be maintained. %However, existing methods are limited to a radius around the prior training data. This poses difficulties for planning when start and goal are in disconnected sets.
% % In many practical robotics applications, we have limited data, and exhaustive coverage of the state-control space with data is infeasible. We may need to temporarily go out of distribution to reach a target, or operate entirely out of distribution and accept degraded performance as long as cautious behavior can enable safety.
% % Existing robust methods cannot handle this.
% % Moreover, these methods are computationally slow, relying on kinodynamic sampling-based planners with wasteful reachable set-environment collision checks that are queried while the planner explores in an uninformed fashion towards the goal. Something missing here is an optimization-based formulation and an error metric that can actively guided the trajectory design. Faster methods can enable real-time replanning, which can also enable online adaptation of model error bounds as new data becomes available, yielding less conservative results and more informed planning.
% % They are also slow. This prevents real-time replanning and adaptation of model error bounds given new data that is being observed online.
% % Emerging tools, SLS can help -- optimization that simultaneously constructs reference trajectories, feedback controllers, and corresponding closed-loop reachable sets. Moreover, weighted conformal prediction (WCP) gives a way to construct error sets that cover the true error with high probability, and coverage gap can be bounded, even in out-of-distribution settings, as long as a bound on how rapidly the two distributions diverge as a function of the input domain is given.
% % To address this gap, we propose \methodnospace, a fast optimization-based feedback motion planning algorithm for robust planning with learned dynamics. \method formulates an MPC problem that allows temporary OOD operation while maintaining probabilistic safety guarantees. 

% To address this gap, we propose \methodnospace, a fast optimization-based feedback planner for robust planning with learned dynamics. In particular, \method formulates an MPC problem that allows temporary OOD operation while maintaining probabilistic safety. 
% During such operation, online data updates the model error bounds to encourage more cautious behavior, such as keeping greater distance from obstacles. \method uses the system level synthesis (SLS) framework \cite{anderson2019system} to jointly design reference trajectories, feedback controllers, and closed-loop reachable tubes. 
% %Weighted conformal prediction (WCP) \cite{barber2023conformal} provides calibrated model error bounds covering true errors with high probability and bounded coverage gap, even OOD, provided the divergence between training and test distributions is limited. 
% %Weighted conformal prediction (WCP) \cite{barber2023conformal} yields calibrated, state–control–dependent error bounds with high-probability coverage and a bounded OOD coverage gap under limited distribution shift.
% %These bounds inform SLS tubes, ensuring high-probability safety when tracking learned dynamics. \method can also include active uncertainty reduction to guide planning toward lower-error regions. 
% These bounds inform the SLS tubes to ensure safety. \method can also bias plans toward low-uncertainty regions.
% Our contributions are: %These bounds inform the SLS-computed tubes, ensuring the closed-loop trajectories remain within them when tracking plans computed with the learned dynamics, with high probability. \method can also include active uncertainty reduction in the SLS cost, guiding planning toward lower-error regions. Our contributions are:

% % To address this gap, we propose \methodnospace, a fast optimization-based feedback motion planning algorithm that enables robust planning with learned dynamics. In particular, \method takes a model predictive control (MPC)-based formulation \cite{rawlings2017model} and allows the system to temporarily leave the training domain while maintaining probabilistic safety guarantees. During such out-of-distribution operation, \method uses data collected online to update its model error bounds to promote more cautious behavior, such as maintaining a larger distance from obstacles. Specifically, \method leverages the system-level synthesis (SLS) \cite{anderson2019system} framework to formulate an optimization problem that simultaneously designs reference trajectories, feedback controllers, and corresponding closed-loop reachable tubes. Weighted conformal prediction (WCP) \cite{barber2023conformal} is used to estimate calibrated model error bounds that cover the true model error with high probability, with a bounded coverage gap even in out-of-distribution settings, provided a bound on how rapidly the training and test distributions diverge across the input domain. These model error bounds then inform the reachable tubes computed by the SLS optimization, ensuring that the closed-loop system trajectories remains within them with high probability when tracking a plan generated with the learned dynamics with an SLS-based feedback controller. Furthermore, \method can incorporate active uncertainty reduction into the SLS cost, guiding the planner toward lower-error regions of the state space. Specifically, our contributions are:
% % To address this gap, we propose \method, a near-real-time optimization-based feedback motion planning algorithm which plans robustly with learned dynamics. In particular, \method enables the system to (temporarily) leave the training domain, while maintaining probabilistic safety guarantees. When doing so, \method can update its model error bounds to encourage more cautious behavior (i.e., steering away from obstacles). Specifically, \method leverages the system level synthesis (SLS) framework to formulate an optimization problem that simultaneously constructs reference trajectories, feedback controllers, and corresponding closed-loop reachable sets. We leverage weighted conformal prediction (WCP) to estimate model error bounds that cover the true model error with high probability, with bounded coverage gap even in out-of-distribution settings, as long as a bound on how rapidly the training and test distributions diverge as a function of the input domain is given. We then use these model error bounds to inform the reachable sets computed by the SLS optimization, bounding the closed-loop behavior of the system when planning with the learned dynamics with high probability. Moreover, we can design active uncertainty-reduction in the SLS cost to actively guide the planner towards lower-error regions of the space.
\looseness-1Provably-safe planning under unknown dynamics is crucial for real-world robotics. This poses difficulties for classical planners, which assume precise models; in practice, model error and uncertainty are unavoidable.
Data-driven planners can learn dynamics from data but may exploit model inaccuracies, leading to unsafe behavior. 
This motivates uncertainty-constrained planners that quantify prediction error between true and learned dynamics to ensure robust, safe execution \citep{DBLP:journals/ral/KnuthCOB21, DBLP:conf/icra/KnuthCRM23, wasiela2024learning}. 
These methods estimate spatially-varying model error bounds to compute reachable tubes, which are used to constrain closed-loop behavior within a ``trusted domain" where the model is accurate. This domain is typically selected as a bounded region around the training data, preventing generated plans from straying far from the data. However, in practice, data scarcity often forces robots to operate temporarily out of distribution (OOD), tolerating degraded performance while maintaining safety.
Moreover, these planners are slow and do not exploit gradient information to efficiently optimize closed-loop behavior that actively reduces model uncertainty.

\looseness-1To address these gaps, we propose \methodnospace, a fast optimization-based robust feedback motion planner that uses the learned dynamics to jointly plan open-loop nominal trajectories and closed-loop tracking controllers via robust model predictive control (MPC). Our method uses weighted conformal prediction (CP) to bound the model error and system level synthesis (SLS) \citep{anderson2019system} to compute reachable sets, which are used to constrain the true closed-loop system to be safe with high probability, enabling robust temporary OOD operation. \method updates the model error bounds with online data to encourage cautious behavior when OOD, such as staying further from obstacles. 
\method also uses gradient data from a model error predictor to efficiently guide plans toward low-error regions, enabling near-real-time replanning. 
Our contributions are:

\vspace{-8pt}
\begin{itemize}
    \itemsep-0.35em
    \item A fast optimization-based MPC method with learned dynamics ensuring finite-sample, probabilistic safety of the true system using state-control-dependent CP error bounds in the MPC.%A fast optimization-based MPC with learned dynamics that ensures finite-sample, probabilistic safety of the true system via SLS and CP, integrating state–control–dependent CP error bounds directly in the MPC.
    \item Theoretical analysis for high-probability safety and bounding the OOD coverage gap.
    \item Numerical validation on a simulated nonlinear 4D car and 12D quadcopter, improving safety and prediction accuracy even in OOD scenarios over uncertainty-unware MPC baselines.
\end{itemize}
% \begin{itemize}
%     \item Fast, optimization-based framework for robust planning with learned dynamics, with probabilistic safety guarantees, using SLS + CP% framework for robust planning with learned dynamics
%     \item Spatially varying CP-based model error bounds that hold with high probability and with bounded coverage gap
%     \item Theoretical analysis, quantified probabilistic error coverage
%     \item Experimental validation in simulation on nonlinear dynamics, a 4D car and 12D quadcopter, outperforming baselines in safety, prediction accuracy
% \end{itemize}
% Experiments:
% \begin{itemize}
%     \item Car (4D)
%     \item Quadcopter (6D)
%     \item Rocket (17D?)
% \end{itemize}

\vspace{-25pt}
\section{Related Work}
\vspace{-4pt}

\paragraph{Predictive Control}
\looseness-1Learning-based MPC unifies disturbance estimation with planning \citep{hewing2020learning}, but typically lacks state-dependent error bounds and cannot steer plans towards low-error regions. Robust \citep{rawlings2017model} and tube MPC \citep{mayne2005robust} ensure constraint satisfaction under disturbance, but in typical formulations, directly enforcing constraints on the closed-loop trajectory is nonconvex, leading to conservative overapproximations \citep{singh2018robust}. In contrast, SLS \citep{anderson2019system,chen2024robust}, or disturbance-feedback MPC \citep{GOULART2006523}, provides a convex parameterization of closed-loop responses for linear time-varying (LTV) systems, enabling constraint satisfaction under disturbance \citep{bartos2025stochastic}. SLS also extends to data-driven settings \citep{xue2021data,furieri2022neural} and nonlinear dynamics \citep{leeman2025robust_TAC}, where the problem is decomposed into an optimization over nominal trajectories and closed-loop LTV error dynamics. In this work, we inform SLS with state-dependent, CP-calibrated error bounds to enable safe planning toward low-error regions.

\vspace{-8pt}
\paragraph{Planning with Learned Dynamics} 
\looseness-1Many model-based reinforcement learning methods plan with learned models but lack safety guarantees \citep{DBLP:journals/ftml/MoerlandBPJ23}. Some methods estimate reachable tubes for safe exploration \citep{berkenkamp2016safe}, but assume a controller is given. The most closely related method \citep{DBLP:conf/cdc/ChouOB21, chou2022safe} safely plans with learned dynamics by constraining invariant tubes to lie in a ``trusted domain" near training data. This demands dense data coverage; in contrast, we enable the robot to move beyond the training data while assuring probabilistic safety. Subsequent model error-aware planners \citep{DBLP:conf/corl/SuhCDYGT23} improve planning speed but lack formal guarantees. Also closely related is \cite{marques2024quantifyingaleatoricepistemicdynamics}, which bounds model error via CP but cannot extrapolate OOD as it uses \textit{exchangeable} CP. Other methods aim to improve dynamics predictions in OOD scenarios \citep{DBLP:journals/corr/abs-2403-12245} but do not focus on planning.

\vspace{-6pt}
\paragraph{Conformal Prediction (CP)}
\looseness-1CP \citep{vovk2005algorithmic} is a distribution-free method for constructing prediction regions from data from black-box models, making it well-suited for safe trajectory prediction and planning \citep{DBLP:journals/ral/LindemannCSP23, DBLP:conf/l4dc/DixitLWCPB23, DBLP:conf/nips/SunJQNKS23, DBLP:conf/cdc/MuthaliSDLRFT23, muenprasitivej2025probabilistically}. The most closely related methods are \cite{DBLP:conf/l4dc/ContrerasSS25}, which triggers a safe fallback policy in OOD scenarios for MPC, and \cite{DBLP:conf/cdc/CheeHP23, DBLP:conf/l4dc/CheeSHP24}, which use WCP to address time-series correlations in MPC for tracking fixed trajectories. However, these methods do not use state-dependent uncertainty, and form prediction sets using a fixed spatially-invariant norm ball. Thus, the MPC cannot distinguish between state-space regions where the model is accurate and those where it is inaccurate, thereby degrading performance. In contrast, our method designs trajectories that are valid OOD and actively encourages uncertainty reduction, steering the MPC toward regions of lower model error, using a learned state-dependent error covariance matrix.

% \paragraph{Predictive Control}

% SLS \citep{anderson2019system,chen2024robust}, equivalent to disturbance-feedback MPC \citep{GOULART2006523}, provides a convex parameterization of closed-loop responses for linear time-varying (LTV) systems, enabling joint optimization of nominal trajectories and feedback. Online feedback optimization reduces conservativeness and improves performance under bounded or stochastic disturbances \cite{bartos2025stochastic}. Extensions to nonlinear dynamics decompose control into a nominal trajectory and an LTV error system, allowing joint optimization of affine feedback and linearization-error bounds \cite{leeman2025robust_TAC}. SLS is also effective in data-driven settings, where learned closed-loop responses enable robust control with limited models \citep{xue2021data,furieri2022neural}.
% Broadly, robust MPC ensures constraint satisfaction under disturbance \cite{rawlings2017model}. Tube MPC \cite{mayne2005robust} uses invariant sets around nominal trajectories. Learning-based MPC methods combine uncertainty estimation with planning \cite{hewing2020learning}, but typically assume simple error models and lack state-dependent bounds.
% % System Level Synthesis (SLS) \cite{anderson2019system,chen2024robust}, equivalent to disturbance-feedback MPC \cite{GOULART2006523}, provides a convex parameterization of closed-loop responses for linear time-varying systems, enabling joint optimization of nominal trajectories and feedback. By optimizing feedback online, SLS reduces conservatism and improves performance under bounded or stochastic disturbances \cite{bartos2025stochastic}. Extensions to nonlinear dynamics decompose control into a nominal trajectory and an LTV error system, allowing joint optimization of affine feedback and linearization-error bounds for robust constraint satisfaction \cite{leeman2025robust_TAC}. SLS is also effective in data-driven settings, where learned closed-loop responses enable robust control with limited models \cite{leeman2023predictive,xue2021data,furieri2022neural,dean2019safely}. 
% % More broadly, robust MPC ensures constraint satisfaction under model mismatch and disturbances \cite{kouvaritakis2016model,rawlings2017model}. Tube-based MPC \cite{langson2004robust,mayne2005robust} uses invariant sets around nominal trajectories, while stochastic MPC \cite{cannon2012stochastic,mesbah2016stochastic} enforces chance constraints. Learning-based MPC methods combine uncertainty estimation with online planning \cite{hewing2020learning,aswani2013provably}, but often rely on Gaussian or Lipschitz error models and lack state-dependent bounds.
% % System Level Synthesis (SLS) \cite{anderson2019system,chen2024robust}, equivalently disturbance-feedback MPC \cite{GOULART2006523}, provides a convex parameterization of closed-loop responses for linear time-varying systems and enables joint optimization of the nominal trajectory and feedback within an MPC loop. SLS avoids compounding over-approximations by optimizing feedback online, which typically reduces conservatism and enhances performance in the presence of bounded or stochastic disturbances \cite{bartos2025stochastic}. Building on this foundation, SLS has been extended to nonlinear dynamics by decomposing control into a nominal trajectory and an LTV error system, enabling joint optimization of affine feedback and linearization-error bounds for robust constraint satisfaction \cite{leeman2025robust_TAC}. Finally, SLS has also proved effective in data-driven settings \cite{leeman2023predictive}, where closed-loop responses are learned from measurements to enable robust control with limited models \cite{xue2021data,furieri2022neural,dean2019safely}.

% % More broadly, robust MPC has been developed to guarantee constraint satisfaction under model mismatch and bounded disturbances \citep{kouvaritakis2016model,rawlings2017model}. Tube-based MPC \citep{langson2004robust,mayne2005robust} uses invariant sets around nominal trajectories, while stochastic MPC \citep{cannon2012stochastic,mesbah2016stochastic} enforces chance constraints using probabilistic models of uncertainty. Recent works on learning-based MPC couple uncertainty estimation with online planning \citep{hewing2020learning,aswani2013provably}, but typically rely on Gaussian or Lipschitz error models and lack state-dependent bounds. 

% \vspace{-8pt}
% \paragraph{Planning with Learned Dynamics} 
% Most model-based RL methods perform long-horizon planning with learned models but lack safety guarantees \cite{DBLP:journals/ftml/MoerlandBPJ23}. Some methods estimate reachable tubes for safe exploration \cite{berkenkamp2016safe}, but these either assume a fixed controller or restrict planning near training trajectories. The most closely related methods \cite{DBLP:conf/cdc/ChouOB21} safely plan with learned dynamics using invariant tubes constrained to lie in a ``trusted domain" near training data, demanding dense coverage. In contrast, we allow the robot to move beyond the training data while retaining probabilistic safety guarantees. Later work \cite{DBLP:conf/corl/SuhCDYGT23} improves speed via diffusion-guided optimization but lacks formal guarantees. Also closely related is \cite{marques2024quantifyingaleatoricepistemicdynamics}, which leverages CP for model error bounding, but unlike our method, cannot extrapolate OOD as it uses exchangeable CP. Other methods aim to improve dynamics model performance in OOD scenarios \cite{DBLP:journals/corr/abs-2403-12245} but do not focus on planning. %Methods leveraging conformal prediction \cite{marques2024quantifyingaleatoricepistemicdynamics} or improving OOD model performance \cite{DBLP:journals/corr/abs-2403-12245} do not support safe extrapolation or provide explicit model error bounds.
% % Most existing model-based RL methods perform long-horizon planning with learned models but lack safety or reachability guarantees. For example, \cite{ichter2019robot} plans in learned latent spaces, and \cite{guzzi2020path} estimates the confidence that a controller can move between states to guide planning. Some emerging methods aim to provide safety assurances. One approach learns stability certificates: \cite{DBLP:conf/nips/KolterM19} and \cite{DBLP:journals/corr/abs-2008-05952} learn Lyapunov functions from data, but applying a single Lyapunov function to general point-to-point planning is challenging. Other methods use Gaussian processes (GPs) for reachable tube estimation \cite{koller2018learning} and safe exploration \cite{akametalu2014reachability, berkenkamp2016safe}, but these assume a given controller; in contrast, we design a feedback controller along with the motion plan. \cite{DBLP:journals/corr/abs-2002-01587} learns trajectory-tracking tubes, yet plans must remain close to the training trajectories, which restricts planning flexibility. Our method allows deviation from the training domain. The most closely related works \cite{DBLP:journals/ral/KnuthCOB21, DBLP:conf/cdc/ChouOB21, DBLP:conf/icra/KnuthCRM23} plan safely with learned dynamics by deriving invariant tubes around a plan within a “trusted domain” near the training data. A key limitation is that the robot must stay within this trusted domain for the guarantees to hold, requiring dense coverage of the state/control space. Our method relaxes this assumption, enabling limited extrapolation beyond the training data. Follow-up work replaces the sampling-based planners of \cite{DBLP:conf/cdc/ChouOB21} with diffusion-guided optimization \cite{DBLP:conf/corl/SuhCDYGT23} to improve speed, but they do not perform reachability analysis for formal guarantees. Also related is \cite{marques2024quantifyingaleatoricepistemicdynamics}, which leverages conformal prediction for model error bounding, but unlike our method, it cannot extrapolate outside the training domain. Other methods improve dynamics model performance in out-of-distribution scenarios \cite{DBLP:journals/corr/abs-2403-12245} but do not explicitly provide model error bounds.

% % Most existing model-based RL methods perform long-horizon planning with learned models without safety or reachability guarantees, e.g., \cite{ichter2019robot} which plans in learned latent spaces and
% % \cite{guzzi2020path}, which estimates the confidence that a controller can move between states to guide planning. Methods are emerging that aim to provide safety assurances. One class of methods learns stability certificates: \cite{DBLP:conf/nips/KolterM19},\cite{DBLP:journals/corr/abs-2008-05952} learn Lyapunov functions from data; however, it is difficult to apply a single Lyapunov function for general point-to-point motion planning. In \cite{DBLP:journals/corr/abs-2005-00611}, the method derives a model-free stability certificate, but does not consider control design. Other methods use Gaussian processes (GPs) for reachable tube estimation \cite{koller2018learning} and safe exploration \cite{akametalu2014reachability, berkenkamp2016safe}, but assume that a controller is given; we design a feedback controller along with the motion plan. \cite{DBLP:journals/corr/abs-2002-01587} learns trajectory tracking tubes; however, plans must remain near full training \textit{trajectories} to be accurate, which can heavily restrict planning; in contrast, our method can deviate away from the training domain. Finally, most relevant are \cite{lipschitz_ral, ChouOB21, DBLP:conf/icra/KnuthCRM23}, which plan safely with learned dynamics by deriving an invariant tube around a plan inside a ``trusted domain" near the training data. A key restriction of \cite{lipschitz_ral, ChouOB21, DBLP:conf/icra/KnuthCRM23} is that the robot must remain within the ``trusted domain" for the guarantees to hold, require dense data coverage in the state/control space to enable practical planning; our method relaxes this assumption, enabling limited extrapolation far from training data. To improve speeds, the robust planner is replaced with a diffusion-guided optimization planner \cite{DBLP:conf/corl/SuhCDYGT23}, but this does not perform reachability analysis to provide formal guarantees. Also closely related is \cite{marques2024quantifyingaleatoricepistemicdynamics}, which leverages conformal prediction for model error bounding, but unlike our method, cannot extrapolate out of the training domain as it uses exchangeable CP. Other methods aim to improve dynamics model performance in out-of-distribution scenarios \cite{DBLP:journals/corr/abs-2403-12245} but do not focus on planning.

% \vspace{-6pt}
% \paragraph{Conformal Prediction (CP)}
% CP \citep{vovk2005algorithmic} is a distribution-free method for constructing prediction regions from data from black-box models, making it well-suited for safe trajectory prediction and planning \citep{DBLP:journals/ral/LindemannCSP23, DBLP:conf/l4dc/DixitLWCPB23, DBLP:conf/nips/SunJQNKS23}. Prior work has applied CP to guide motion planning \citep{DBLP:conf/cdc/MuthaliSDLRFT23}. The most closely related methods are \cite{DBLP:conf/l4dc/ContrerasSS25}, which triggers a safe fallback policy in OOD scenarios for MPC, and \cite{DBLP:conf/l4dc/CheeSHP24,DBLP:conf/cdc/CheeHP23}, which use WCP to address time-series correlations in MPC for tracking fixed trajectories. However, these methods do not use state-dependent uncertainty. In contrast, our method designs trajectories online in OOD settings and explicitly encourages uncertainty reduction, steering the MPC toward regions of lower model error.

% Conformal prediction \citep{vovk2005algorithmic} is a black-box, distribution-free approach for constructing prediction regions from observed data, making it suitable for safe trajectory prediction and planning \citep{DBLP:journals/ijrr/LuoZKISSP24, DBLP:journals/ral/LindemannCSP23, DBLP:conf/l4dc/DixitLWCPB23, DBLP:conf/nips/SunJQNKS23, muenprasitivej2025probabilisticallysafebipedalnavigationuncertain}. Prior work has used CP to verify autonomous systems \citep{cp_baseline} and guide motion planning \citep{DBLP:journals/ral/LindemannCSP23, DBLP:conf/cdc/MuthaliSDLRFT23}. Most closely related are \cite{DBLP:conf/l4dc/ContrerasSS25} which uses CP to trigger a safe fallback policy when controlling in OOD scenarios and \cite{DBLP:conf/l4dc/CheeSHP24,DBLP:conf/cdc/CheeHP23}, which leverage weighted CP in MPC for tracking a fixed set of trajectories without active uncertainty reduction; in contrast, our method designs trajectories online and explicitly promotes uncertainty reduction in the cost to guide the MPC to regions with low model error. \glen{anutam, check and expand on this} In our method, we use CP to design an MPC strategy for planning with learned dynamics that enables probabilistic safety guarantees for the true closed-loop system. Unlike prior work, we leverage SLS and weighted CP to enable computation of robust closed-loop reachable sets under dynamics error in OOD scenarios.

\vspace{-10pt}
\section{Preliminaries and Problem Statement}\label{sec:preliminaries}
\vspace{-4pt}

Let $f: \mathcal{X} \times \mathcal{U} \rightarrow \mathcal{X}$ be the true unknown deterministic nonlinear discrete-time dynamics,

\vspace{-9pt}
\begin{equation}\label{eq:dynamics}
    x_{k+1} = f(x_k, u_k) = \hat f(x_k, u_k) + \varepsilon(x_k, u_k),
\end{equation}
\vspace{-13pt}

\looseness-1\noindent where $\mathcal{X} \subseteq \mathbb{R}^{n_x}$ is the state space and $\mathcal{U} \subseteq \mathbb{R}^{n_u}$ is the control
space. Let $\hat f: \mathcal{X} \times \mathcal{U} \rightarrow \mathcal{X}$ denote an approximate learned dynamics model, where $\varepsilon: \mathcal{X} \times \mathcal{U} \rightarrow \mathcal{X}$ is the dynamics error. We model $\hat f$ with a neural network (NN), though our method is agnostic to model class. The goal is to reach a target set $\mathcal{X}_\textrm{f} \subseteq \mathcal{X}$ while satisfying pointwise-in-time state and control constraints $(x_k, u_k) \in \mathcal{F}$, where

\vspace{-18pt}
\begin{equation}\label{eq:constraints}
    \mathcal{F} := \{(x, u) \in \mathbb{R}^{n_x + n_u} \mid g_i(x,u) + b_i \le 0,\quad \forall i\in\{1,\ldots,n_c\}\text{,\quad for }g_i:\mathbb{R}^{n_x + n_u} \to \mathbb{R}\}.
\end{equation}
\vspace{-16pt}

\noindent Let $\mathcal{B}^m \subseteq \mathbb{R}^m $ be the unit 2-norm ball. Let the Minkowski sum be $A \oplus B := \{a + b\mid a \in A, b \in B\}$.% Our method relies heavily on CP and SLS, which we overview next.

\vspace{-5pt}
% \paragraph{Conformal Prediction (New)}
% We use conformal prediction (CP) to convert a raw non-conformity score into a calibrated threshold that yields $(1-\alpha)$ coverage \citep{vovk2005algorithmic}. Let $\calib:={(x_i,u_i,f(x_i,u_i))}_{i=1}^n$ and choose a score $s:\mathcal{X}\times\mathcal{U}\times\mathcal{X}\to\mathbb{R}{\ge0}$. Set $s_i:=s(x_i,u_i,f(x_i,u_i))$. Calibration proceeds in three steps:
% \begin{itemize}
%     \item[(i)] form the weighted empirical score law with optimized $\tilde w$
% \begin{equation}\label{eq:empirical_distribution}\textstyle
% S:=\sum_{i=1}^n \tilde w_i,\delta_{s_i}+\tilde w_{\textrm{test}},\delta_{+\infty},\qquad
% \tilde w_{\textrm{test}}+\sum_{i=1}^n \tilde w_i=1,
% \end{equation}
% \item[(ii)] take the calibrated cutoff $q_{1-\alpha}:=Q_{1-\alpha}(S)$, and
% \item[(iii)]  define the prediction set at $(x_{\textrm{test}},u_{\textrm{test}})$ as
% \begin{equation}\label{eq:prediction_set}
%     \mathcal{C}(x_\textrm{test}, u_\textrm{test}; q_{1-\alpha}):= \{x \in \mathcal{X} \mid s(x_\textrm{test}, u_\textrm{test}, x) \leq q_{1-\alpha}\}.
% \end{equation}
% \end{itemize}

% Classical CP uses uniform weights and relies on exchangeability; under spatio-temporal shift this may fail. We therefore use weighted CP (WCP) \citep{barber2023conformal}, choosing importance weights $\tilde w_i$ (normalized as above) to emphasize calibration points whose score distribution matches the test query. Let $S^i$ denote the distribution obtained by replacing $s_i$ with the unknown test score. WCP guarantees
% \vspace{-8pt}
% \begin{equation}\label{eq:cp_guarantee}
% \Pr\big[f(x_{\textrm{test}},u_{\textrm{test}})\in \mathcal{C}(x_{\textrm{test}},u_{\textrm{test}};q_{1-\alpha})\big]
% \ge 1-\alpha-\sum_{i=1}^n \tilde w_i\tvd(S^{i},S),
% \end{equation}
% %\vspace{-10pt}
% so the weighted total-variation term quantifies the coverage gap due to shift (it vanishes under exchangeability). In practice, the weights can be selected to trade off a small gap $\sum_i \tilde w_i,\tvd(S^{i},S)$ and compact sets $\mathcal{C}(x_{\textrm{test}},u_{\textrm{test}};q_{1-\alpha})$ (e.g., via Gaussian kernel), yielding uncertainty sets suitable for robust MPC with learned dynamics.

\paragraph{Conformal Prediction (CP)}
We use weighted CP to convert raw non-conformity scores, quantifying a mismatch between between $\hat f(x,u)$ and $f(x,u)$, into a calibrated threshold that yields $(1-\alpha)$ coverage \citep{vovk2005algorithmic}. Let $\calib:={(x_i,u_i,f(x_i,u_i))}_{i=1}^n$ and choose a nonconformity score function  $s:\mathcal{X}\times\mathcal{U}\times\mathcal{X}\to\mathbb{R}_{\ge0}$. Set $s_i:=s(x_i,u_i,f(x_i,u_i))$. Calibration proceeds in three steps. We
(i) form the weighted empirical distribution $\hat{S}$ with weights, $\{\tilde{w}_{1}, \ldots, \tilde{w}_{n}, \tilde w_\textrm{test}\} \subseteq [0,1]^{n+1}$,

\vspace{-8pt}
\begin{equation}\label{eq:empirical_distribution}\textstyle
\hat{S}:=\sum_{i=1}^n \tilde w_i\delta_{s_i}+\tilde w_{\textrm{test}}\delta_{+\infty},\quad \textrm{s.t. }
\tilde w_{\textrm{test}}+\sum_{i=1}^n \tilde w_i=1,
\end{equation}
\vspace{-13pt}

\noindent where $\delta_{s_i}$ is the Dirac delta centered at $s_i$, (ii) calibrate the threshold $q_{1-\alpha}:=Q_{1-\alpha}(\hat{S})$ with the $1-\alpha$ quantile $Q_{1-\alpha}(\cdot)$, and (iii) define the conformal prediction set (CP set) at $(x_{\textrm{test}},u_{\textrm{test}})$ as

\vspace{-10pt}
\begin{equation}\label{eq:prediction_set}
    \mathcal{C}(x_\textrm{test}, u_\textrm{test}; q_{1-\alpha}):= \{x \in \mathcal{X} \mid s(x_\textrm{test}, u_\textrm{test}, x) \leq q_{1-\alpha}\}.
\end{equation}
\vspace{-15pt}
% \begin{itemize}
%     \item[(i)] form the weighted empirical distribution with weights, $\{\tilde{w}_{1}, \dots, \tilde{w}_{n}, \tilde w_\textrm{test}\} \subseteq [0,1]^{n+1}$ 
% \begin{equation}\label{eq:empirical_distribution}\textstyle
% S:=\sum_{i=1}^n \tilde w_i,\delta_{s_i}+\tilde w_{\textrm{test}},\delta_{+\infty},\quad \textrm{s.t. }
% \tilde w_{\textrm{test}}+\sum_{i=1}^n \tilde w_i=1,
% \end{equation}
% \item[(ii)] calibrate the threshold $q_{1-\alpha}:=Q_{1-\alpha}(S)$ with the $1-\alpha$ quantile $Q_{1-\alpha}(\cdot)$, and
% \item[(iii)]  define the conformal prediction set (CP set) at $(x_{\textrm{test}},u_{\textrm{test}})$ as
% \begin{equation}\label{eq:prediction_set}
%     \mathcal{C}(x_\textrm{test}, u_\textrm{test}; q_{1-\alpha}):= \{x \in \mathcal{X} \mid s(x_\textrm{test}, u_\textrm{test}, x) \leq q_{1-\alpha}\}.
% \end{equation}
% \end{itemize}

\noindent\looseness-1 Traditional (non-exchangeable) CP uses uniform weights, $\tilde w_i = \tilde w_\text{test} = \frac{1}{n+1}$,  and assumes exchangeability, which may fail under spatio-temporal distribution shifts; WCP does not assume this. 
% We select importance weights $\tilde w_i$ (normalized as above) to emphasize calibration points whose score distribution matches the test query. 
We denote $(s_1, s_2, \ldots, s_\textrm{test}) \sim S^\textrm{test}$; moreover, let $S^i$ be the joint distribution obtained by replacing $s_i$ with the unknown test score, i.e., $(s_1, s_2, \ldots, s_{i-1}, s_\textrm{test}, s_{i+1}, \ldots, s_n) \sim S^i$. WCP guarantees
% Traditional (non-exchangeable) CP uses uniform weights, $\tilde w_i = \tilde w_\text{test} = \frac{1}{n+1}$,  and relies on exchangeability which may fail under spatio-temporal shifts. Thus, we use weighted CP (WCP). 
% % We select importance weights $\tilde w_i$ (normalized as above) to emphasize calibration points whose score distribution matches the test query. 
% Then, let $S^i$ denote the distribution obtained by replacing $s_i$ with the unknown test score. WCP guarantees
\vspace{-5pt}
\begin{equation}\label{eq:cp_guarantee}
\Pr\big[f(x_{\textrm{test}},u_{\textrm{test}})\in \mathcal{C}(x_{\textrm{test}},u_{\textrm{test}};q_{1-\alpha})\big]
\ge 1-\alpha-\textstyle\sum_{i=1}^n \tilde w_i\tvd(S^{i},S^\textrm{test}),
\end{equation}
\vspace{-17pt}

\looseness-1\noindent where $d_\textrm{TV}(\cdot)$ is the total-variation distance and $\textstyle\sum_{i=1}^n \tilde w_i\tvd(S^{i},S^\textrm{test})$ quantifies the coverage gap due to distribution shift (it vanishes under exchangeability). In practice, the weights can be selected to trade off coverage gap size $\sum_i \tilde w_i\tvd(S^{i},S^\textrm{test})$ and CP set size $|\mathcal{C}(x_{\textrm{test}},u_{\textrm{test}};q_{1-\alpha})|$, yielding uncertainty sets suitable for robust MPC with learned dynamics. %(e.g., via Gaussian kernel)
WCP is essential for planning with learned dynamics to account for (1) distribution shifts between training and execution data, (2) state/control correlations during execution, and (3) visits to OOD regions of the state/control space.
% so the weighted total-variation term quantifies the coverage gap due to shift (it vanishes under exchangeability). In practice, the weights can be selected to trade off a small gap $\sum_i \tilde w_i,\tvd(S^{i},S)$ and compact sets $\mathcal{C}(x_{\textrm{test}},u_{\textrm{test}};q_{1-\alpha})$ (e.g., via Gaussian kernel), yielding uncertainty sets suitable for robust MPC with learned dynamics.
% For planning with learned dynamics, WCP is essential to account for (1) distribution shifts between training data and correlated states/controls during execution and (2) visits to OOD regions of the state/control space



% \paragraph{Conformal Prediction (old)} CP is an uncertainty quantification (UQ) framework that produces prediction sets covering the true outcome with probability $1-\alpha$, where $\alpha \in (0,1)$ is a user-selected miscoverage rate \citep{vovk2005algorithmic}. 
% %Traditional CP assumes statistical exchangeability, meaning the data is identically distributed and all permutations of the data ordering are equally likely. This limits its use when spatio-temporal distribution drifts occur, e.g., when a robot plans far from its training data. To address this, we use non-exchangeable weighted conformal prediction (WCP), which relaxes the exchangeability assumption while retaining similar coverage guarantees \citep{barber2023conformal}.
% Classical CP assumes exchangeability (i.i.d. data with random ordering), which fails under spatio-temporal drift, e.g., planning far from training data. We instead use weighted conformal prediction (WCP), which relaxes this assumption while preserving coverage guarantees \citep{barber2023conformal}.

% %This precludes its use in situations where spatio-temporal distributional drifts occur, e.g., when a robot plans to move far from the training data of its dynamics model. Thus, we consider non-exchangeable (or weighted) conformal prediction (WCP) that relaxes the exchangeability assumption while maintaining similar coverage guarantees as traditional CP \citep{barber2023conformal}.

% Given a calibration set $\calib := {(s_i)}_{i=1}^n$ with $s_i := s(x_i, u_i, f(x_i, u_i))$, CP/WCP use a non-conformity score $s:\mathcal{X}\times\mathcal{U}\times\mathcal{X}\to\mathbb{R}_{\ge 0}$ to quantify mismatch between $(x_i,u_i)$ and $f(x_i,u_i)$. Larger $s$ means greater disagreement; by itself this score is uncalibrated and carries no coverage guarantee.
% %Given a  $n$ points calibration dataset $\calib := \{(x_i, u_i, f(x_i, u_i)\}_{i=1}^n$, CP and WCP leverages a non-conformity score function $s: \mathcal{X}\times\mathcal{U}\times\mathcal{X} \to \mathbb{R}_{\geq 0}$ to compute an \textit{uncalibrated} notion of uncertainty, i.e., lacking probabilistic coverage guarantees, where a larger non-conformity score indicates greater disagreement between the input $(x_i, u_i)$ and the target $f(x_i, u_i)$. 
% Hence, to calibrate coverage, we form a probability measure

% \vspace{-12pt}
% \begin{equation}\label{eq:empirical_distribution}\textstyle
%     S := \sum_{i=1}^n\tilde{w}_{i}\delta_{s_i} + \tilde{w}_\textrm{test} \delta_{+\infty},
% \end{equation}
% %\vspace{-18pt}
% where $\tilde{w}_{i}$ are weights normalized such that $\tilde w_\textrm{test} + \sum_{i=1}^{n}\tilde{w}_{i} = 1$ and $\delta_{s_i}$ refers to the Dirac delta function centered at $s_i$. All weights $\{\tilde{w}_{i}\}_{i=1}^n$ and $\tilde w_\textrm{test}$ lie within the interval $[0, 1]$. We denote the $(1-\alpha)$-quantile of the probability measure $S$ as $q_{1-\alpha} := Q_{1-\alpha}(S)$ and %The weights ${w}_{i}$ prior to normalization satisfy ${w}_{i} \in [0,1]$, for all $i \in\{1, \ldots, n+1\}$, with ${w}_{n+1}=1$ -- a convention set in \cite{barber2023conformal}.%$\infty$ representing the unknown inference point, where
% define the prediction set at test point $(x_\textrm{test}, u_\textrm{test})$ as

% \vspace{-10pt}
% \begin{equation}\label{eq:prediction_set}
%     \mathcal{C}(x_\textrm{test}, u_\textrm{test}; q_{1-\alpha}):= \{x \in \mathcal{X} \mid s(x_\textrm{test}, u_\textrm{test}, x) \leq q_{1-\alpha}\}.
% \end{equation}
% \vspace{-15pt}

% We also define $S^i$ as the distribution of the $n$-tuple $(s_1, \dots, s_{i-1},s_\textrm{test}, s_{i+1}, \dots s_n)$, when we replace the $i^\mathrm{th}$ calibration point with the test point. Then, the WCP method ensures that 

% \vspace{-10pt}
% \begin{equation}\label{eq:cp_guarantee}
%     \textstyle\Pr[f(x_\textrm{test}, u_\textrm{test}) \in \mathcal{C}(x_\textrm{test}, u_\textrm{test};q_{1-\alpha})] \geq 1 - \alpha -\sum_{i=1}^{n}\tilde{w}_i\,\tvd(S^{i}, S),
% \end{equation}
% \vspace{-14pt}

% \noindent where $\tvd(S^{i}, S)$ is the total variation distance between distribution $S^{i}$ and $S$.
% The weights $\tilde{w}_i$ can be tuned to jointly minimize the coverage gap $\sum_{i=1}^{n}\tilde{w}_i\,\tvd(S_{i}, S)$ and the size of the prediction set $\mathcal{C}(x_\textrm{test}, u_\textrm{test};q_{1-\alpha})$. Intuitively, by weighting calibration scores more similarly distributed to the test non-conformity score, we can reduce the coverage gap. Under the standard exchangeability assumption made in traditional CP, $\tvd(S^{i}, S)=0$ and the gap is eliminated. WCP is crucial for planning with learned dynamics, as it accounts for (1) distribution shifts between training data and correlated states/controls during execution and (2) visits to OOD regions of the state/control space.%The WCP framework is essential for planning with learned dynamics by accounting for 1) distribution shifts between training data and correlated states/controls during execution and 2) visits to OOD regions of the state/control space. %\glen{one sentence on why we can tune these weights arbitrarily and maintain probabilistic guarantees. one sentence on how this is different to regular CP. one sentence on the intuition of the WCP framework}%find a Pareto-optimal point to 

\vspace{-8pt}
\paragraph{SLS} SLS \citep{anderson2019system} is a robust control framework for uncertain LTV systems

\vspace{-10pt}
\begin{equation}\label{eq:ltv_dynamics}
x_{k+1} = A_kx_{k} + B_ku_{k} + E_k\xi_k    
\end{equation}
\vspace{-15pt}

\looseness-1\noindent where $A_k\in\mathbb{R}^{n_x \times n_x}$, $B_k\in\mathbb{R}^{n_x \times n_u}$, and $E_k \in \mathbb{R}^{n_x \times n_x}$, and $\xi_k \in \mathcal{B}^{n_x}$ is a disturbance. SLS can be used to jointly optimize a nominal trajectory, a tracking controller, and a closed-loop reachable tube overapproximation computed by propagating worst-case disturbances over time. In particular, we denote a nominal length-$T$ state trajectory $\textbf{z} := \{z_1, \ldots, z_T\} \in \mathbb{R}^{n_x T}$ and control trajectory $\textbf{v} := \{v_1, \ldots, v_T\} \in \mathbb{R}^{n_u T}$ satisfying the nominal LTV dynamics in \eqref{eq:nominal_ltv}. To track the nominal $(\textbf{z}, \textbf{v})$, SLS defines a disturbance-feedback controller \eqref{eq:ctrl_sls} and closed-loop state response \eqref{eq:state_sls}, where $\Pu_{k,j} \in \mathbb{R}^{n_u \times n_x}$ and $\Px_{k,j} \in \mathbb{R}^{n_x \times n_x}$ define the closed-loop response of \eqref{eq:ltv_dynamics} as an affine function of $\{\xi_k\}$,
%models uncertain disturbances \citep{anderson2019system}. 
\vspace{-30pt}
\begin{multicols}{3}
    \begin{equation}\label{eq:nominal_ltv}
    z_{k+1} = A_k z_k + B_k v_k,
\end{equation}

    \begin{equation}\label{eq:ctrl_sls}
    u_k = v_k + \textstyle\sum_{j=1}^{k-1}\Pu_{k,j}\xi_j,
  \end{equation}

  \begin{equation}\label{eq:state_sls}
    x_k = z_k + \textstyle\sum_{j=1}^{k-1}\Px_{k,j}\xi_j.
  \end{equation}
\end{multicols}
\vspace{-5pt}

\looseness-1\noindent To ensure that \eqref{eq:ctrl_sls}-\eqref{eq:state_sls} represent the closed-loop state/control response, $\Px_{k,j}$ and $\Pu_{k,j}$ must satisfy 

\vspace{-8pt}
\begin{equation}\label{eq:sls_affine}
    \Px_{k+1,j} = A_k\Px_{k,j} + B_k\Pu_{k,j},\qquad \Px_{k+1,k} = E_k.
\end{equation}
\vspace{-13pt}

\noindent With SLS, we can guarantee that the closed-loop trajectories of \eqref{eq:state_sls} and \eqref{eq:ctrl_sls} remain in the reachable tubes 
%Under these conditions SLS guarantees, that the trajectory $(x_k, u_k) \in \mathcal{R}_k = \mathcal{R}^{x}_k \times \mathcal{R}^{u}_k$ remains in tubes 
$\mathcal{R} := \{\mathcal{R}_k:= (\mathcal{R}_k^x, \mathcal{R}_k^u)\}_{k=1}^T$, i.e., for all $k \in \{1, \ldots, T\}$, $x_k \in \mathcal{R}_k^x$ and $u_k \in \mathcal{R}_k^u$, where

\vspace{-5pt}
\begin{equation}\label{eq:sls_tubes}\textstyle
    \mathcal{R}^{x}_k := (\bigoplus_{j=0}^{k-1}\Px_{k,j}\mathcal{B}^{n_x})\oplus \{z_k\}, \quad \mathcal{R}^{u}_k := (\bigoplus_{j=0}^{k-1}\Pu_{k,j}\mathcal{B}^{n_x})\oplus \{v_k\}.
\end{equation}
\vspace{-15pt}

% Those tubes are constrained within the state and input constraints for each time steps, i.e., $(\mathcal{R}^x_k \times \mathcal{R}^u_k) \subseteq \mathcal{X} \times \mathcal{U}$. 
% \noindent Minimizing a surrogate volume of the reachable tubes can be achieved through a cost term $\tilde H_0(\Px,\Pu)$, discussed in App. \ref{app:probs}, where $\Px$ and $\Pu$ collects all $\Px_{k,j}$ and $\Pu_{k,j}$ respectively.
\noindent SLS has also been extended to control nonlinear dynamics \citep{leeman2025robust_TAC,zhan2025robustly} with robust satisfaction of nonlinear constraints, both via optimized linearization error overbounds. For nonlinear dynamics $\hat f$ of the form \eqref{eq:dynamics}, we can obtain an LTV approximation
using the Jacobian linearizations of $\hat f$ around $(\textbf{z}, \textbf{v})$: $A_k:= \bm{\nabla}_x \hat{f}\left(z_k, v_k\right),~\text{and}~B_k:=\bm{\nabla}_u \hat{f}\left(z_k, v_k\right)$.
To robustly satisfy nonlinear constraints of the form \eqref{eq:constraints}, we can enforce the tightened constraints \eqref{eq:nonlinear_constraint}, obtained by linearizing $g_i$, for all $i \in \{1,\ldots, n_c\}$, around the nominal trajectory \eqref{eq:nominal_ltv}: 
\begin{equation}\label{eq:nonlinear_constraint}\textstyle
\sum_{j=0}^{k-1}\big\Vert\big[\nabla_xg_i(z_k,v_k)\Px_{k,j}, \nabla_ug_i(z_k,v_k)\Pu_{k,j}\big]\big\Vert_{2} + g_i(z_k,v_k) + b_i +g_i^\textrm{lin}(z_k,v_k) \leq 0,
\end{equation}
where $g_i^\textrm{lin}(z_k,v_k)$ is a linearization error bound term (details are omitted for brevity; see \cite[Eq. (22)]{zhan2025robustly}). This ensures that the closed-loop dynamics satisfy \eqref{eq:constraints} for all $k \in \{1,\ldots,T\}$.

\paragraph{Problem Statement}

% We assume that we are given an offline dataset $\mathcal{D} := \{(x_i, u_i, f(x_i, u_i))\}_{i=1}^{N}$, where each $(x_i, u_i)$ is drawn from a joint distribution $Y$. We partition $\mathcal{D}$ into a training dataset $\mathcal{D}_\textrm{train} := \{(x_i, u_i, f(x_i, u_i))\}_{i=1}^{N_\textrm{train}}$ and calibration dataset $\mathcal{D}_\textrm{calib} := \{(x_i, u_i, f(x_i, u_i)\}_{i=1}^{N_\textrm{calib}}$. We will leverage $\mathcal{D}_\textrm{train}$ to train our models -- the learned dynamics model $\hat f$ \eqref{eq:dynamics} and a heuristic model error predictor $\mathbf{\Sigma}$, described in Sec. \ref{sec:methods}. Our goal is to 1) estimate calibrated error bounds on the learned dynamics $\hat f$ that hold with high probability and to 2) use these bounds within a robust MPC framework, with a prediction horizon of $T$ steps, to guarantee robust constraint satisfaction with high probability. Specifically, we solve: 
We assume that we are given an offline dataset $\mathcal{D} := \{(x_i, u_i, f(x_i, u_i))\}_{i=1}^{N}$, where each $(x_i, u_i)$ is drawn from a joint distribution $Y$. We partition $\mathcal{D}$ into a \textit{training dataset} $\mathcal{D}_\textrm{train} := \{(x_i, u_i, f(x_i, u_i))\}_{i=1}^{N_\textrm{train}}$, to train the learned dynamics $\hat f$ \eqref{eq:dynamics} and a model error predictor $\mathbf{\Sigma}: \mathcal{X} \times \mathcal{U} \rightarrow \mathbb{R}^{n_x \times n_x}$, described in Sec. \ref{sec:methods}, and \textit{calibration dataset} $\mathcal{D}_\textrm{calib} := \{(x_i, u_i, f(x_i, u_i)\}_{i=1}^{N_\textrm{calib}}$. Our goal is to 1) estimate calibrated error bounds on the learned dynamics $\hat f$ that hold with high probability and to 2) use these bounds within a robust MPC framework, with a prediction horizon of $T$ steps, to guarantee robust constraint satisfaction with high probability. Specifically, we solve: 

%an approximate learned dynamics model $\hat f$ \eqref{eq:dynamics} and error predictor $\mathbf{\Sigma}$ which outputs covariance matrices that capture the 

% Here, we formulate a robust model predictive control, explicitly accounting for the data distribution used to generate the residual dynamics and its uncertain bounds. Those bounds are leveraged to construct the disturbance propagation and corresponding constraint tightening.

% In this regard, we seek to solve the following problems:\\
% \noindent\textit{\underline{Problem 1: Learning and Calibration.}} Using $\mathcal{D}_\textrm{train}$, learn an approximate dynamics model $\hat f: \mathcal{X} \times \mathcal{U} \rightarrow \mathcal{X}$ and a heuristic uncertainty model $\bm \Sigma: \mathcal{X} \times \mathcal{U} \rightarrow \mathbb{R}^{n_x \times n_x}$ which predicts $n_x \times n_x$ covariance matrices that estimate the size of the dynamics uncertainty. Using $\hat f$, $\mathbf{\Sigma}$, and $\calib$, provide a calibrated high-confidence bound $E: \mathcal{X} \times \mathcal{U} \rightarrow 2^\mathcal{X}$ on the residual dynamics $\varepsilon: \mathcal{X}\times \mathcal{U}\rightarrow \mathbb{R}^{n_x}$, such that $\textrm{Pr}[\bigcap_{k=1}^T \big(\varepsilon(x_k,u_k) \in E(z_k,u_k)\big)] \ge 1-\alpha-\sum_{k=1}^{T-1}\sigma_k$, for nominal state $(z_k, u_k)$ and for all reachable state/controls $(x_k,u_k) \in \mathcal{R}_k:= (\mathcal{R}_k^x, \mathcal{R}_k^u)$, as defined in \eqref{eq:sls_tubes}. 
% \noindent\textit{\underline{Problem 1: Learning and Calibration.}} Using $\mathcal{D}_\textrm{train}$, learn an approximate dynamics model $\hat f: \mathcal{X} \times \mathcal{U} \rightarrow \mathcal{X}$. Using $\hat f$ and $\calib$, provide a calibrated high-confidence bound $E: \mathcal{X} \times \mathcal{U} \rightarrow 2^\mathcal{X}$ on the residual dynamics $\varepsilon: \mathcal{X}\times \mathcal{U}\rightarrow \mathbb{R}^{n_x}$, such that $\textrm{Pr}[\bigcap_{k=1}^T \big(\varepsilon(x_k,u_k) \in E(z_k,u_k)\big)] \ge 1-\alpha-\sum_{k=1}^{T-1}\sigma_k$, for nominal state $(z_k, u_k)$ and for all reachable state/controls $(x_k,u_k) \in \mathcal{R}_k:= (\mathcal{R}_k^x, \mathcal{R}_k^u)$, as defined in \eqref{eq:sls_tubes}. Here, $\sigma_k$ refers to the coverage gap \eqref{eq:cp_guarantee} at timestep $k$.
\noindent\textit{\underline{Problem 1: Learning and Calibration.}} Using $\mathcal{D}_\textrm{train}$, learn an approximate dynamics model $\hat f: \mathcal{X} \times \mathcal{U} \rightarrow \mathcal{X}$. Using $\hat f$ and $\calib$, provide a calibrated high-confidence bound $\mathcal{E}: \mathcal{X} \times \mathcal{U} \rightarrow 2^\mathcal{X}$ on the dynamics error $\varepsilon: \mathcal{X}\times \mathcal{U}\rightarrow \mathbb{R}^{n_x}$, such that $\textrm{Pr}[\bigcap_{k=1}^T \big(\varepsilon(x_k,u_k) \in \mathcal{E}(z_k,v_k)\big)] \ge 1-\tilde\alpha$, for nominal state $(z_k, v_k)$ and for all reachable state/controls $(x_k,u_k) \in \mathcal{R}_k:= (\mathcal{R}_k^x, \mathcal{R}_k^u)$, as defined in \eqref{eq:sls_tubes}. Here, $\alpha \in(0,1)$ is user-specified and $\tilde \alpha \in [\alpha, 1)$ absorbs the coverage gap in \eqref{eq:cp_guarantee}.

%Intuitively, we wish to ensure that $E(z_k, u_k)$ contains the true error $e(x_k, u_k)$ for all $k \in \{1,\ldots,T\}$ and for all state/controls that may be reached, as defined by \eqref{eq:sls_tubes}, with probability at least $1-\alpha$.% $E(x_k,u_k)$ such that $\varepsilon(x_k,u_k) \in E(x_k,u_k)$ using $\calib$.\\

\looseness-1\noindent\textit{\underline{Problem 2: Planning.}} Using the learned model $\hat f$ and error bound $\mathcal{E}$ of Prob. 1, solve a robust MPC problem for $\hat f$ ensuring that the constraints \eqref{eq:constraints} are satisfied for the true closed-loop dynamics (i.e., under $f$ in \eqref{eq:dynamics}) with probability at least $1-\tilde \alpha$, jointly over $T$ prediction steps, for a single MPC solve, and with probability at least $1-\hat \alpha$ for $R$ MPC solves (i.e., $R-1$ replanning steps), where $\hat \alpha \in (\tilde \alpha, 1)$. 
% \red{Antoine: Is $\tilde a = \tilde \alpha$?}
% \red{Not clear if you consider a per step probability of the joint probability. (The per step probability is a really strong assumption that typically doesn't hold.)}
%to perform forward reachability analysis across the MPC horizon and develop constraint tightenings for each step in the horizon. 
% Below we discuss non-exchangeable Conformal Prediction (CP) and System-Level-Synthesis (SLS), which we use to address problems 1 \& 2 in Section \ref{sec:methods}.

\vspace{-10pt}
\section{Method}\label{sec:methods}
We outline our method, \methodnospace, which uses WCP for high-probability state-control-dependent model error bounding (Sec. \ref{sec:cp_err_bound}) and integrates it into an uncertainty-constrained SLS-based robust MPC optimization (Sec. \ref{sec:sls_plus_cp_bound}). Finally, we discuss practical implementation (Sec. \ref{sec:scp}). %we consider Problems 1 \& 2 in Sections \ref{sec:cp_err_bound} and \ref{sec:sls_plus_cp_bound} respectively. 

\vspace{-12pt}
\subsection{Model Training and Model Error Bounding via Conformal Prediction}\label{sec:cp_err_bound}
\looseness-1To tackle Prob. 1, we first train the dynamics $\hat f$ \eqref{eq:dynamics} on the training dataset $\mathcal{D}_\textrm{train}$ by minimizing a mean-squared-error loss $\textstyle\sum_{k=1}^{N_\textrm{train}} \Vert f(x_k, u_k) - \hat f(x_k, u_k)\Vert_2$. Then, for any nominal state $z \in \mathcal{X}$ and control input $v \in \mathcal{U}$, we seek to bound the error of the learned dynamics model $\Vert \hat{f}(z, v)-f(z,v)\Vert_2 $. To encode spatially-varying error, we train a heuristic uncertainty model $\covariance: \mathcal{X}\times \mathcal{U} \to \mathbb{R}^{n_x \times n_x}$, an NN, with a multivariate Gaussian negative-log-likelihood (MGNLL) loss \citep{dasgupta2007line}:% omitting the constant $n_x \ln(2\pi) $,
\begin{equation}\textstyle\small
\mathcal{L}(\covariance; f, \hat{f}, x, u) = \frac{1}{2} \left( (f(x,u) -\hat{f}(x,u))\T \covariance(x,u)^{-1} (f(x,u) -\hat{f}(x,u)) + \ln(\det(\covariance(x,u)))\right).
\end{equation}
% \red{Antoine: Here, you implicitly assume that the residual model follows a Gaussian distribution? And you need this assumption for your results to hold, isn't it? Or is this part of the "conformalization" step?}
\looseness-1MGNLL is selected to train $\mathbf{\Sigma}$ to output an uncalibrated covariance matrix which locally approximates the \textit{shape} and \textit{magnitude} of $\hat{f}$'s error dispersion at a given state-control pair. WCP's calibration procedure reconciles the inexact gaussian assumption (of MGNLL) made of the error dispersion by scaling the covariance matrix.   
To capture the model error at timestep $k$ for a nominal ($z_k, v_k)$, we adapt the non-conformity score of \cite{messoudi2022ellipsoidal}:
\begin{equation}\label{eq:non_conform}\small
s_{i,k} := s_k(x_i, u_i, f(x_i, u_i)) = \sqrt{\big(f(x_i, u_i) - \hat{f}(x_i, u_i)\big )\T\covariance(z_k,v_k)^{-1}\big(f(x_i, u_i) - \hat{f}(x_i, u_i)\big)}.
\end{equation}
Using \eqref{eq:non_conform}, we form the $(1-\alpha_k)$-confidence CP set for some nominal $(z_k, v_k)$, $\mathcal{C}(z_k, v_k; {q_{1-\alpha_k}})$, as
\begin{equation}\label{eq:ellip_prediction_set}\small
    \mathcal{C}(z_k,v_k; {q_{1-\alpha_k}}):= \{x \in \mathcal{X} \mid \sqrt{\big(x - \hat{f}(z_k,v_k)\big )\T\covariance(z_k,v_k)^{-1}\big(x - \hat{f}(z_k,v_k)\big)} \leq q_{1-\alpha_k}(z_k, v_k)\},
\end{equation}
defining $q_{1-\alpha_k}(z_k,v_k):=Q_{1-\alpha_k}(\hat{S}_k)$, where $\hat{S}_k := \sum_{i=1}^n \tilde w^k_i\delta_{s_{i,k}}+\tilde w^k_{\textrm{test}}\delta_{+\infty}$ is formed using a set of normalized weights $\tilde W^k := \{\tilde{w}_{1}^{k}, \dots, \tilde{w}_{N_\textrm{calib}}^{k}, \tilde w_\textrm{test}^k\} \subseteq [0,1]^{N_\textrm{calib}+1}$ tailored to the nominal state/control pair $(z_k,v_k)$. 
Intuitively, $\mathcal{C}(z_k,v_k; {q_{1-\alpha_k}})$ is an ellipsoid centered at the predicted next state $\hat{f}(z_k, v_k)$, with a radius given by the quantile of the weighted non-conformity scores from \eqref{eq:non_conform}. 

To construct the normalized weights $\tilde W^k$ tailored for a nominal state/control pair $(z_k,v_k)$, we first define a set of unnormalized weights $W^k = \{{w}_{1}^{k}, \dots, {w}_{N_\textrm{calib}}^{k}, w_\textrm{test}^k=1\} \subseteq [0,1]^{N_\textrm{calib}+1} $, where
\begin{equation}\label{eq:rho_weights}
    w^k_i = \rho^{||(z_k,v_k) - (x_i, u_i)||_2},
\end{equation}
\looseness-1where $\rho \in (0, 1]$, for all $(x_i, u_i)\in \calib$. The form of \eqref{eq:rho_weights} encodes the assumption that model error is state/control-dependent and changes smoothly, such that nearby pairs have similar errors. We normalize the weights $w^k$ as $\tilde w^k_i = \frac{w^k_i}{1 + \sum_{j=1}^{n}w^k_j}$ to be compatible with the WCP framework (Sec. \ref{sec:preliminaries}). 

\looseness-1Denote $(s_{1,k}, s_{2,k}, \ldots s_{i-1,k}, s_{k,k}, s_{i+1,k}, \ldots, s_{N_\textrm{calib},k}, s_{i,k}) \sim S^{i,k}$ and $(s_{1,k}, \ldots, s_{N_\textrm{calib},k}, s_{k,k}) \sim S^{k,k}$ (i.e., $S^{i,k}$ exchanges $s_{k,k}$ with $s_{i,k}$), where $s_{k,k} = s_k(z_k,v_k,f(z_k,v_k))$ is generally unknown as $f(z_k,v_k)$ may not be in $\calib$ for some candidate $(z_k, v_k)$. Via the WCP framework, \eqref{eq:cp_guarantee} ensures

\vspace{-8pt}
\begin{equation}\label{eq:cp_guarantee_f}\textstyle
    \textrm{Pr}[{f}(z_k, v_k) \in \mathcal{C}(z_{k}, v_{k};{q_{1-\alpha_k}})] \ge 1 - \alpha_k -\sum_{i=1}^{n}\tilde{w}_i^k\,\tvd(S^{i,k}, S^{k,k}).
\end{equation}
\vspace{-12pt}

\noindent Starting from \eqref{eq:cp_guarantee_f}, we see that the dynamics error is in an origin centered ellipsoid, $\varepsilon(z_k,v_k) = \big(f(z_k,v_k) - \hat f(z_k, v_k) \big) \in \mathcal{C}(z_k, v_k; q_{1-\alpha_k}) \oplus \{-\hat f(z_k, v_k)\}$, implying that it satisfies $\Pr[\varepsilon(z_k,v_k) \in \mathcal{E}(z_k,v_k)] \geq 1 - \alpha_k -\sum_{i=1}^{n}\tilde{w}_i^k\,\tvd(S^{i,k}, S^{k,k}) $, where we set $\mathcal{E}(z_k, v_k) := \mathcal{C}(z_k, v_k; q_{1-\alpha_k}) \oplus \{-\hat f(z_k, v_k)\}$. In essence, by conformalizing the heuristic uncertainty model $\mathbf{\Sigma}$, we obtain a calibrated spatially-dependent model error bound $\mathcal{E}(z_k, v_k)$ that holds with high probability. 
% Specifically, the weights $w^k_i$ are dependent on the nominal point $(z_k,v_k)$, 
Equivalently, using the Cholesky factor $L(z_k,v_k)$ of $\mathbf{\Sigma}(z_k,v_k)$, i.e., $\mathbf{\Sigma}(z_k,v_k) = L(z_k,v_k)L(z_k,v_k)^\top$, \eqref{eq:ellip_prediction_set}-\eqref{eq:cp_guarantee_f}, and selecting a miscoverage rate $\alpha_k \in (0,1)$ tailored to timestep $k$, we can state that
\begin{equation}\label{eq:ball_form}\textstyle
\textrm{Pr}[\varepsilon(z_k,v_k) \in \underbrace{V(z_k,v_k)\mathcal{B}^{n_x}}_{:=\mathcal{E}(z_k,v_k)}] \ge 1 - \alpha_k -\sum_{i=1}^{N_\textrm{calib}}\tilde{w}_i^k\,\tvd(S^{i,k}, S^{k,k}),
\end{equation}\vspace{-22pt}

\begin{equation}\label{eq:ball_form_vzk}\textstyle
\text{where}\quad V(z_k,v_k) := {q_{1-\alpha_k}(z_k,v_k)L(z_k,v_k)}.
\end{equation}
% To encode state-dependence into the prediction, we use a weighting scheme where the weights $w^k_i$ are dependent on the nominal point $(z_k,v_k)$, $w^k_i = \rho^{||(z_k,v_k) - (x_i, u_i)||}$ for $(x_i, u_i)\in \calib$. The corresponding normalized weights are $\tilde w^k_i = \frac{w^k_i}{1 + \sum_{j=1}^{n}w^k_j}$. To encode state-dependent uncertainty we train a neural-network $\bm\Sigma: \mathcal{X}\times \mathcal{U} \to \mathcal{X}$ with multivariate gaussian negative-log-likelihood (MGNLL) loss omitting the constant $n_x \ln(2\pi) $,
% \begin{equation}\textstyle
% \mathcal{L}(\boldsymbol{\Sigma}; f(x,u), \hat{f}(x,u)) = \frac{1}{2} \left( (f(x,u) -\hat{f}(x,u))^T \boldsymbol{\Sigma}^{-1} (f(x,u) -\hat{f}(x,u)) + \ln(\det(\boldsymbol{\Sigma}))\right),
% \end{equation}
% to provide covariance matrices indicating relative and correlated uncertainty for each dimension in the state-space \citet{dasgupta2007line}. Then we compute non-conformity scores, $s_i$ using Equation \ref{eq:non_conform}, replacing $\Sigma$ with $\Sigma(z_k,v_k)$. 
% While \eqref{eq:ball_form} bounds model error at $(z_k, v_k)$, during execution the system may be perturbed within the reachable tube $\mathcal{R}_k$. Since $(x,u) \neq (z_k,v_k)$ in general, it is important to quantify coverage of $E(z_k, v_k)$ from any $(x,u) \in \mathcal{R}_k$. The ellipsoids $E(z_k,v_k)$ are computed a priori at the nominal point, so we can augment the coverage gap:
% \[
% \Pr[\varepsilon(x,u) \in E(z_k,v_k)] > 1-\alpha - 2\sum_{i=1}^{n} \tilde{w}_i d_{TV}(S_i, \tilde{S}_k) - \gamma(\mathcal{R}_k),
% \]
% where $s_i \sim S_i$, $\tilde{S}_k$ is the distribution of $s_i(z_k, v_k, f(z_k, v_k))$, and $\gamma(\mathcal{R}_k)$ is defined in Thm.~\ref{thm:tube_gap}.

% Finally, the probability that the error holds for all timesteps in the trajectory is:

% While \eqref{eq:ball_form} is sufficient for bounding model error when starting exactly from $(z_k, v_k)$, during execution, our system may be perturbed within the reachable tube $\mathcal{R}_k$ at timestep $k$. Thus, since we are not guaranteed that $(x,u) = (z_k,v_k)$, it is essential to quantify the coverage of $E(z_k, v_k)$ starting from any $(x,u) \in \mathcal{R}_k$. Since the ellipsoids $E(z_k,v_k)$ are computed a-priori for the nominal point $(z_k, v_k)$, we can augment the coverage gap
\looseness-1\noindent While \eqref{eq:ball_form}-\eqref{eq:ball_form_vzk} bounds model error at $(z_k, v_k)$, the system may be perturbed during execution within the reachable tube $\mathcal{R}_k$. Thus, at every timestep $k$, we must quantify the coverage of $\mathcal{E}(z_k, v_k)$ from any $(x_k,u_k) \in \mathcal{R}_k$. 
% This can be achieved by inflating the coverage gap with an additional term $\gamma(\mathcal{R}_k)$ (defined in Thm. \ref{thm:tube_gap}), yielding $\Pr[\varepsilon(x_k,u_k) \in \mathcal{E}(z_k,v_k)] \ge 1-\alpha_k- 2\sum_{i=1}^{N_\textrm{calib}} \tilde{w}_i^k \tvd(S_{i,k}, S_{k,k}) - \gamma(\mathcal{R}_k)$, where $S_{i,k}$ and ${S}_{k,k}$ are distributions such that $s_{i,k} \sim S_{i,k}$, $s_{k,k} \sim {S}_{k,k}$. 
This can be achieved by inflating the coverage gap with an additional term, yielding $\Pr[\varepsilon(x_k,u_k) \in \mathcal{E}(z_k,v_k)] \ge 1-\alpha_k- 2\sum_{i=1}^{N_\textrm{calib}} \tilde{w}_i^k \tvd(S_{i,k}, S_{k,k}) - \gamma(\mathcal{R}_k)$, where $\gamma(\mathcal{R}_k)$ is defined in Thm. \ref{thm:tube_gap}. Here, $S_{i,k}$ and ${S}_{k,k}$ are distributions such that $s_{i,k} \sim S_{i,k}$, $s_{k,k} \sim {S}_{k,k}$. 
Specifically, Thm. \ref{thm:total_coverage} derives the probability that $\varepsilon(x_k,u_k)\in\mathcal{E}(z_k,v_k)$ for all $(x_k,u_k) \in \mathcal{R}_k$, for all $k \in \{1,\ldots,T\}$, providing our solution for Prob. 1 (proofs in App. \ref{app:proofs}):

\vspace{-4pt}
\begin{theorem}\label{thm:total_coverage}
    \looseness-1Given a length-$T$ trajectory of the true dynamics $f$ \eqref{eq:dynamics} $\tau_f:= \{(x_k, u_k)\}_{k=1}^{T}$ such that $(x_k, u_k) \in \mathcal{X} \times \mathcal{U}$ for all $k \in \{1, \ldots, T\}$, desired miscoverage $\{\alpha_k\}_{k=1}^T$, total miscoverage $\sigma_k = \alpha_k + 2\sum_{i=1}^{N_\textrm{calib}} \tilde{w}_i^k \tvd(S_{i,k}, S_{k,k})+ \gamma({\mathcal{R}_k})$, and independence between  $\calib$ and $(x_k,u_k)$, we have:

    \vspace{-10pt}
    \begin{equation}\textstyle
        \Pr[\bigcap_{k=1}^{T}~\big(\varepsilon(x_k,u_k) \in \mathcal{E}(z_k,v_k)\big)] \ge 1- \sum_{k=1}^{T}\sigma_k := 1-\tilde\alpha.
    \end{equation}
\vspace{-20pt}
    
\end{theorem}
Overall, at every timestep, we predict an ellipsoid $\mathcal{E}(z_k,v_k)$ containing the one-step error $\varepsilon(x_k,u_k)$ with high probability. Next, we will use this data-driven, spatially-dependent bound $\mathcal{E}(z_k, v_k)$ to inform SLS-based robust MPC, shrinking tubes where data is dense and expanding them if OOD.

% \red{Antoine:I think we need a tighter connection between SLS and CP here.\\At each time step, we have an ellipsoid, for the one-step model error. We over-approximate by the scaled $E_k$, thereby integrating data-driven, state-dependent uncertainty directly into the SLS tightenings above. This preserves tractability while shrinking tubes where data are dense and expanding them in OOD regions.}
\vspace{-10pt}
\subsection{Informing SLS-based Planning with CP Model Error Bounds}\label{sec:sls_plus_cp_bound}

\looseness-1Using our conformalized model error bounds, we construct an uncertain nonlinear system $z_{k+1} = \hat{f}(z_k,v_k) + V(z_k,v_k)\xi_k$,
where $\xi_k \in \mathcal{B}^{n_x}$ and $ V(z_k,v_k)\xi_k$ represent the unknown, bounded model error at timestep $k$. 
Let $\Px$ and $\Pu$ collect all $\Px_{k,j}$ and $\Pu_{k,j}$. Then we adapt \cite{leeman2025robust_TAC} to write a nonlinear program (NLP) in \eqref{eq:nonlinear_sls} to optimize for $\Z = \{z_1,\dots,z_{T}\}$, $\V = \{v_1,\dots,v_{T-1}\}$, $\Px$, and $\Pu$, given an initial state $\bar{x}_0 \in \mathcal{X}$: 
% \antoine{We have $\Z = \{z_1,\dots,z_{T}\}$, but the optimization problem uses $z_{T+1}$.}
% \red{Antoine: The optimization problem does not use $v_T$. So when is it included in the variables? Are you assuming that $\mathcal{F} = \mathcal{X}_f$? It is not so intuitive to have the input in the terminal set, especially in an MPC paper.}
\begin{subequations}\small
\label{eq:nonlinear_sls}
\begin{align}
\min_{\substack{\mathbf{\Phi}_{\mathbf{x}},\mathbf{\Phi}_{\mathrm{u}},\mathbf{z}, \mathbf{v}}}\quad &  J(\mathbf{z},\mathbf{v}) +  \tilde H_0(\Px, \Pu) + J_{\mathcal{X}_\textrm{f}}(\textbf{z}, \textbf{v}),\label{eq:nonlinear_objective}\\[-6pt]
\text { s.t. }\quad\quad &  {z}_{k+1}= \hat{f}(z_k,v_k),~ {z}_1 = \bar x_0,\quad\quad \forall k = 1,\ldots, T-1,\label{eq:nonlinear_nominal}\\
&\Px_{k+1,j} =A_k  \Px_{k,j} + B_k \Pu_{k,j},\ \quad \forall j = 1,\ldots, T-1,\quad \forall k = j+1,\ldots,T-1,\label{eq:SLP_nonlinear}\\
&\Px_{j+1,j} = V(z_j, v_j),\quad \forall j = 1,\ldots, T-1,\label{eq:SLP_ic_nonlinear}\\%~~\text{(Equation \ref{eq:ball_form})}
& \textstyle\sum_{j=1}^{k}\big\Vert[\nabla_xg_{i}(z_k,v_k)\Px_{k,j}, \nabla_ug_{i}(z_k,v_k)\Pu_{k,j}]\big\Vert_{2} + g_{i}(z_k,v_k) + b_{i} + g_i^\textrm{lin}(z_k,v_k)  \leq 0,\label{eq:SLC_nonlinear}\\
&  \quad\quad \forall i =1, \ldots, n_c,\quad \forall k= 1,\ldots,T. \nonumber
\end{align}
\end{subequations}
% \vspace{-20pt}
\noindent Here, \eqref{eq:nonlinear_objective} penalizes a cost $J(\mathbf{z},\mathbf{v})$ on the nominal trajectory,  $J_{\mathcal{X}_\textrm{f}}(\textbf{z}, \textbf{v})$ penalizes reaching the terminal set $\mathcal{X}_\textrm{f}$, and $\tilde H_0(\Px, \Pu)$ penalizes the size of the reachable tubes (see \eqref{eq:regulizer}, \eqref{eq:lqr}, \eqref{eq:lqr_terminal}). 
% \red{Antoine: If the goal is the reach the terminal set $\mathcal{X}_\textrm{f}$, then it should be used somewhere in the optimization problem. At the moment, only $g_i$ are used that are from the definition of $\mathcal{F}$.}
Since $\hat f$ is more accurate in the training domain, we modify $J(\mathbf{z},\mathbf{v})$ to include an active error reduction cost $J_\mathrm{active}(\mathbf{z},\mathbf{v})$ (see App.~\ref{app:uncert_reduc}) minimizing the distance of $(\mathbf{z},\mathbf{v})$ to in-distribution calibration points. Constraint \eqref{eq:nonlinear_nominal} ensures nominal dynamic feasibility, while \eqref{eq:SLP_nonlinear} and \eqref{eq:SLP_ic_nonlinear} enforce SLS conditions \eqref{eq:sls_affine} for the LTV approximation obtained by linearizing $\hat f$ around $(\mathbf{z},\mathbf{v})$, with model error bound $V(z_j,v_j)$ from \eqref{eq:ball_form_vzk}. Constraint \eqref{eq:SLC_nonlinear} enforces robust state-control satisfaction using bounds on the closed-loop response derived by informing SLS with CP. Since SLS guarantees an overapproximation of the true reachable tube, (1) the closed-loop trajectory executing $(\mathbf{z}, \mathbf{v}, \Px, \Pu)$ over the $T$-step horizon satisfies \eqref{eq:constraints} with high probability, and (2) safety is maintained with high, though diminishing, probability over iterative replanning steps. The following result addresses Prob. 2:
% Here, the objective function \eqref{eq:nonlinear_objective} penalizes a cost on the nominal trajectory $J(\mathbf{z},\mathbf{v})$ and the size of the reachable tubes $\tilde H_0(\Px, \Pu)$ (as defined in \eqref{eq:regulizer}). Since $\hat f$ will be more accurate in the training domain, we minimize an active error reduction cost $J_\text{active}(\mathbf{z},\mathbf{v})$ which minimizes the distance of $(\textbf{z}, \textbf{v})$ to in-distribution calibration points, which serves as a surrogate of the model error (see App. \ref{app:uncert_reduc} for details). Constraint \eqref{eq:nonlinear_nominal} provides a nominal trajectory that is dynamically-feasible with respect to the learned dynamics model. Constraints \eqref{eq:SLP_nonlinear} and \eqref{eq:SLP_ic_nonlinear} encode the SLS conditions \eqref{eq:sls_affine} for an LTV approximation which arises from linearizing the learned dynamics around the nominal trajectory $(\textbf{z}, \textbf{v})$, where the model error bound $V(z_j, v_j)$ is defined in \eqref{eq:ball_form}. Finally, \eqref{eq:SLC_nonlinear} enforces robust satisfaction of the state-control constraints \eqref{eq:constraints} imposed at every timestep, using disturbance bounds derived from the CP model error bounds \eqref{eq:ball_form}. %Constraint \eqref{eq:SLP_ic_nonlinear} uses the CP model error bound from \eqref{eq:ball_form} to initialize the SLS disturbance , $\Px_{j+1,j}$, at time-step $j$.
% Since SLS deterministically guarantees an overapproximation of the reachable tube given a compact disturbance set, it follows that (1) the closed-loop trajectory obtained by executing the $(\textbf{z}, \textbf{v}, \Px, \Pu)$ from \eqref{eq:nonlinear_sls} over the full $T$-step MPC planning horizon on the true dynamics \eqref{eq:dynamics} satisfies the constraints \eqref{eq:constraints} with high probability, and (2) safety is maintained with high, but diminishing, probability across replanning steps. Formally, we have:

\vspace{-6pt}
\begin{theorem}\label{cor:rollout_prob}
    Define the forecasted nominal trajectory {\small $\tau_{\hat f}^{(q)} := \{(z_k^{(q)}, v_k^{(q)})\}_{k=1}^{T}$} as the $T$-timestep future nominal trajectory predicted by the learned dynamics {\small$\hat f$} at the $q$th MPC solve (i.e., the $(q-1)$st replanning step) in \eqref{eq:nonlinear_nominal}. \textbf{(a)} Then the true closed-loop dynamics \eqref{eq:dynamics} under the controller defined by \eqref{eq:nonlinear_sls} and \eqref{eq:ctrl_sls} satisfies {\small $\Pr[\bigcap_{k=1}^T \big((x_k^{(q)}, u_k^{(q)})\in \mathcal{F}\big)] \geq 1 - \sum_{k=1}^{T}\sigma_k^{(q)}$}, where {\small $\{(x_k^{(q)}, u_k^{(q)})\}_{k=1}^T$} is the realized state-control trajectory if executed for the full $T$ timesteps, and \textbf{(b)} defining $\sigma_1^{(q)}$ as the miscoverage rate for the first timestep at MPC solve iteration $q$, we have that by executing \eqref{eq:nonlinear_sls} over $R$ MPC solves (i.e., $R-1$ re-planning steps), {\small $\Pr[\bigcap_{q=1}^R\big((x_1^{(q)}, u_1^{(q)})\in \mathcal{F}\big)] \geq 1 - \sum_{q=1}^{R}\sigma^{(q)}_1 := 1-\hat \alpha$}. 
\end{theorem}
\vspace{-4pt}
% \red{Antoine: For large R (number of replanning step), the probability of constraint satisfaction is actually 0, right? Let me know if you disagree}
\noindent Overall, by solving \eqref{eq:nonlinear_sls}, we reduce model error by 1) minimizing reachable tube volume and 2) penalizing distance to data. Coupling planning with CP-informed reachability ensures high-probability safety for the true dynamics. Finally, refining the CP error set $\mathcal{C}$ with online data (Alg.~\ref{alg:robust_mpc}) reduces the coverage gap, improving empirical model error coverage.
% Overall, our approach is able to reason about and reduce the accumulation of model error by optimizing an objective that 1) minimizes the volume of the closed-loop reachable tubes and 2) penalizes distance to data. This coupling of planning with CP-informed reachability ensures safety under the true dynamics with high probability. Moreover, by incorporating online data, we continuously refine the CP error set $\mathcal{C}$ (Alg. \ref{alg:robust_mpc}), thereby reducing the coverage gap and improving the empirical coverage of realized model errors.

%Benefits: actively reduce uncertainty via optimizing objective, informing planning with downstream capabilities of the tracking controller, which are influenced by the local model error bound. Enabled by coupling planning and CP-informed closed-loop reachability analysis to ensure safety for the true dynamics with high probability. Moreover, we can react to online data, reducing coverage gap and increase empirical coverage of realized model errors in the CP set $\mathcal{C}$, by augmenting the dataset with observed online data (Alg. \ref{alg:robust_mpc}).

% Since, our learned dynamics model will perform better in the training domain, we also consider the active uncertainty reduction cost which seeks to minimize the distance to (in domain) calibration points. Our cost function and approach is discussed in Appendex \ref{app:uncert_reduc}
% \glen{anutam, talk about the active uncertainty reduction cost}

% Here’s a version reduced to about **70%** of the original length while keeping all key technical details:

% ---

% To efficiently solve the nonconvex NLP \eqref{eq:nonlinear_sls}, we propose \methodnospace, summarized in Alg.~\ref{alg:robust_mpc}. \method leverages sequential convex programming (SCP) \citep{malyuta2022convex}, where each iteration is a second-order cone program (SOCP) (line \ref{alg:lin:single}) obtained by linearizing the learned dynamics (\hat f) around the nominal trajectory ((\mathbf{z}, \mathbf{v})) and convexifying constraint tightenings \eqref{eq:SLC_nonlinear}. 
\vspace{-8pt}

\subsection{Practical Implementation via SCP}\label{sec:scp}
\vspace{-4pt}

\looseness-1To efficiently solve the NLP \eqref{eq:nonlinear_sls}, we propose \method (Alg. \ref{alg:robust_mpc}). \method uses sequential convex programming (SCP) \citep{malyuta2022convex,messerer2021survey}. Each SCP iteration solves a second-order cone program (SOCP) (Alg. \ref{alg:robust_mpc}, line \ref{alg:lin:single}), obtained by linearizing $\hat f$ around the current nominal trajectory $(\textbf{z}, \textbf{v})$ and convexifying the constraint tightenings in \eqref{eq:SLC_nonlinear}, yielding an update $(\Z\l^*, \V\l^*)$ to the nominal trajectory. 
To solve this SOCP efficiently, we use the solver in \cite{leeman2024fast}, which iterates between solving a QP for $(\mathbf{z}, \mathbf{v})$ and Riccati recursions for $(\Px, \Pu)$ (line \ref{alg:lin:fastsls}). For efficiency, we apply a real-time iteration scheme \citep{Gros02012020}, limiting SCP to one iteration per step \citep{leeman2025guaranteed}, and omit linearization error \citep{zhan2025robustly}, which is small in practice (Sec.~\ref{sec:results}). At each step, we execute the first control input $v_1$ (line \ref{alg:lin:execute}), shift the previous solution for initialization (line \ref{alg:lin:shift}), and re-solve via SCP. Moreover, we improve OOD robustness by updating the error set $\mathcal{V}$ with online observations (line \ref{alg:lin:data}). Augmenting the error set updates the quantile $q_{1-\alpha_k}(z_k, v_k)$, which in turn alters constraint \eqref{eq:SLP_ic_nonlinear} via \eqref{eq:ball_form_vzk}, reducing the coverage gap and increasing the likelihood that the closed-loop dynamics remain within the computed reachable tubes.
% Then using the solver in \cite{leeman2024fast}, we solve this SOCP by iterating between solving a quadratic program (QP) to obtain the nominal trajectory $(\textbf{z}, \textbf{v})$ and solving Riccati recursions to obtain the disturbance feedback matrices $(\Px, \Pu)$ (line \ref{alg:lin:fastsls}). For practical computation efficiency, we use a real-time iteration scheme \cite{Gros02012020}, i.e., instead of solving the SCP to full convergence at every step, we limit it to one iteration \citep{leeman2025guaranteed}. We also omit the linearization error from \eqref{eq:nonlinear_sls}, following \cite{zhan2025robustly}, as it is minimal with small time-discretization, as is verified empirically in Sec. \ref{sec:results}. We execute the resulting first control input $v_1$ (line \ref{alg:lin:execute}), shift the previous solution for initialization (line \ref{alg:lin:shift}), and re-solve the SCP at the next timestep. Specific to our CP-based SLS approach, to improve OOD robustness, we update our set of errors $\mathcal{V}$ with errors observed during online execution (line \ref{alg:lin:data}). This update reduces the coverage gap and increases the likelihood that the closed-loop system remains within the computed reachable tubes.

\vspace{-4pt}
\begin{algorithm}[H]\small
\caption{\method}\label{alg:robust_mpc}
\begin{algorithmic}[1]
\Procedure{CP-SLS-MPC}{$\calib, \text{horizon }T, \mathcal{F}, \bar x_0, \text{ terminal set }\mathcal{X}_\f$}
    \State $t \gets 1$ and $(\Z, \V) \gets $ Nominal Solution \Comment{Initialize the problem with start state $\bar{x}_0$}
    \State $\mathcal{V} \gets \{(x_i, u_i, f(x_i,u_i) - \hat f(x_i,u_i))\}_{i=1}^{N_\textrm{calib}},\ \forall (x_i, u_i, f(x_i, u_i))\in\calib$ \Comment{Initialize error set $\mathcal{V}$}%$\mathcal{V} \gets \{(x_i,u_i, \hat f(x_i,u_i)) \in \calib | (x_i, u_i, f(x_i,u_i) - \hat f(x_i,u_i))\}$
    \State \textbf{while} {$x_t \notin \mathcal{X}_\f$} \textbf{do} \Comment{$\mathcal{X}_\f$ is included as a terminal cost \eqref{eq:lqr}, App \ref{app:lqr_app}}
        \State $\quad \bar x_0 \gets x_t$; initialize $(\Z, \V)$ with shifted previous solution \Comment{Initialize with the current state}
        % \State $\quad$Initialize $(\Z, \V)$ with shifted previous solution \Comment{Horizon length $T$}
        \label{alg:lin:shift}\State $\quad$Linearize \eqref{eq:nonlinear_sls} around $(\Z, \V)$; compute $V(z_k, v_k)$ via \eqref{eq:ball_form_vzk}, $\mathcal{V}$; form SOCP
        \label{alg:lin:single}\State $\quad$$(\Z\l^*, \V\l^*)$, $\Px, \Pu \gets$ SOCP solution \Comment{Single iteration of \cite{leeman2024fast}}.
        \label{alg:lin:fastsls}\State $\quad$Update $(\Z, \V) \leftarrow (\Z, \V) + (\Z\l^*, \V\l^*)$.

        \State $\quad$$u_t \gets v_1$; $x_{t+1} \gets f(x_t,u_t) $ \Comment{Closed loop control; state update}
        \label{alg:lin:execute}\State $\quad \mathcal{V} \gets \mathcal{V} \cup \{(x_t,u_t, x_{t+1} - \hat{f}(x_t,u_t))\}$; $t \gets t + 1$ \Comment{Online data update}\label{alg:lin:data}
\EndProcedure
\end{algorithmic}
\end{algorithm}
\vspace{-26pt}

\section{Theoretical Analysis: Bounding the Coverage Gap}\label{sec:theory}
%\vspace{-6pt}

% \paragraph{Bounding Coverage Gap:} 
\looseness-1To bound the coverage gap, we assume non-conformity scores from the calibration set $\calib$ are independent of the score at the nominal point $(z_k,v_k)$, though they may come from different distributions. We also assume a Lipschitz-type bound on the distribution drift between ${S_{i,k}}$ and $S_{k,k}$, i.e.,
% To bound our coverage gap, we assume non-conformity scores from our calibration set $\calib$ are independent of the non-conformity score for nominal point $(z,v)$, but may come from non-identical distributions. We also assume a Lipschitz-type bounded distribution drift between $\tilde{S}$ and $S_i$ i.e.,

\vspace{-10pt}
\begin{equation}\label{eq:lipschitz}
\tvd( S_{i,k},S_{k,k}) \leq  \epsilon ||[z_k\T,v_k\T] - [x_i\T,u_i\T]||_{2},    
\end{equation}
\vspace{-13pt}

\noindent for $\epsilon \in \mathbb{R}, \epsilon > 0$, where $s_i \sim S_i$ \eqref{eq:non_conform} and $\tilde{S}$ is the non-conformity score distribution of $(z,v)$. 
\eqref{eq:lipschitz} is a practical assumption, stating that the dynamics error distribution shifts gradually over the state-control space. 
The next theorem shows a overbound on the calibration error $\tvd(S_{i,k},S_{k,k})$.
\vspace{-8pt}
\begin{theorem}\label{thm:cov_gap_thm} The \textit{Coverage Gap}$\,=\sum_{i=1}^{N_\textrm{calib}}\tilde{w}_i\,\tvd(S^{i,k}, S^{k,k})$ of our conformal predictor satisfies

\vspace{-10pt}
\begin{equation}\textstyle
    \text{Coverage Gap}\le \underbrace{\textstyle\sum_{i=1}^{N_\textrm{calib}} \left( \frac{\rho^{d_i}}{1 + \sum_{j=1}^{N_\textrm{calib}} \rho^{d_j}} \right) \cdot 2\epsilon \cdot d_i}_{\text{Tight Bound}} \le \underbrace{\textstyle 2 \epsilon \left[  \frac{d_1 }{1 - \rho^{d_\text{min}}} + \frac{d_\text{max}\rho^{d_\text{min}}}{(1-\rho^{d_\text{min}})^2}\right]}_{\text{Interpretable Bound}}, 
\label{eq:thm3_bound}\end{equation}
where we sort the distance to points in the calibration set such that $d_i = ||[z_k\T,v_k\T] - [x_i\T,u_i\T]||_{2}$, $d_1 \le d_2 \le \dots \le d_n$, $d_\text{min} = \min_{1 \leq i \leq n-1}\{d_{i+1} - d_i\}$, and $d_\text{max} = \max_{1 \leq i \leq n-1} \{ d_{i+1} - d_i\}$.
\end{theorem}
\vspace{-2pt}

% Theorem \ref{thm:cov_gap_thm} implies that the coverage gap is based on the spread of calibration points in the joint state-control space. For a greater density of points, $d_\text{min}$ would decrease, thus increasing the coverage gap. 
\noindent Thm. \ref{thm:cov_gap_thm} (interpretable bound) shows that the gap depends on state-control calibration point density. Greater point density decreases $d_\text{min}$, increasing the gap; reducing $\rho$ mitigates this, but doing so can cause the $(1-\alpha)$-quantile to diverge. The term $\frac{d_1}{1 - \rho^{d_\text{min}}}$ shows that being closer to the training data reduces the gap, as $d_1$ is the minimal distance between $\calib$ and $(z_k,v_k)$. Further discussion of the intuition and validation of the empirical bounds is in App. \ref{app:theoretical_expansion}. Online augmentation of $\mathcal{V}$ may violate independence (due to time-correlation), but for large $N_\textrm{calib}$, Thm. \ref{thm:cov_gap_thm} suffices. As SLS requires model error to hold for all $(x,u) \in \mathcal{R}_k$, we also provide a coverage bound valid over the reachable tube:
% \red{Antoine: "may slightly violate". I think it does violate independence, so I would remove the word \textbf{slightly}, such that it is clear that Theorem 1 not longer apply.}
% \vspace{-15pt}
% To mitigate this, $\rho$ can be reduced, but its reduction should be limited to ensure the $(1-\alpha)$-quantile of the empirical distribution does not degenerate to $\infty$. Also, the term $\frac{d_1}{1 - \rho^{d_\text{min}}}$ implies that being closer to the training distribution will reduce the coverage gap, as $d_1$ is the smallest distance between $\calib$ and the nominal point $(z,v)$. Finally, by augmenting the dataset $\mathcal{V}$ online, we may encounter minor violations to the independence assumption (though the error distribution may still be uncorrelated between successive points). However, with large $N_\textrm{calib}$, Theorem \ref{thm:cov_gap_thm} is a sufficient model for bounding coverage gap. Since SLS assumes that the model error holds for all potential state-control pairs in the reachable tube $\mathcal{R}$, we also provide a coverage bound that holds for all $(x,u)\in \mathcal{R}$:
\begin{theorem}\label{thm:tube_gap}
Under the same assumption of Theorem \ref{thm:cov_gap_thm}, we have that for all $(x,u) \in \mathcal{R}_k$, the tube around the nominal point $(z_k,v_k)$, 
% \begin{equation}\textstyle
    $\Pr[\varepsilon(x,u) \in \mathcal{E}(z_k,v_k)] \geq 1-\alpha_k- 2\sum_{i=1}^{n} \tilde{w}_i d_{TV}(S_{i,k}, S_{k,k})- \gamma({\mathcal{R}_k})$,
% \end{equation}
for $ \gamma({\mathcal{R}_k}) := 2\hat{\epsilon}M(\mathcal{R}_k) $, where $M(\mathcal{R}_k)$ is the maximum axis length of tube $\mathcal{R}_k$, and $\hat \epsilon \le \epsilon$ is chosen such that \eqref{eq:lipschitz} holds for all $(x, u) \in \mathcal{R}_k$.%, and is less than the global Lipschitz constant in $\epsilon$, Theorem \ref{thm:cov_gap_thm}.
\end{theorem}


Using Thm. \ref{thm:tube_gap}, we find that the coverage guarantees remain close to the desired coverage level $\alpha$, assuming a gradual spatial-drift of the error distribution. Furthermore, the optimization problem \eqref{eq:nonlinear_sls} seeks to minimize the tube volume in the cost function, thereby minimizing $\gamma({\mathcal{R}}_k)$. 

\vspace{-12pt}
\section{Results}\label{sec:results}
% Metrics to report:
% \begin{itemize}
%     \item Goal error (averaged over $\approx$10 runs)
%     \item Average tracking error (averaged over $\approx$10 runs)
%     \item Max tracking error
%     \item Tube sizes for method
%     \item Computation time (and explain why it is slow)
%     \item Dataset sizes (show that the computation time is a function of the dataset size)
% \end{itemize}
\vspace{-3pt}

% \subsection{Baselines and Experiment Setup}
We evaluate our method on robust planning with learned 4D Dubins' (see App.~\ref{app:car_dyn}) and 12D quadcopter dynamics (see App. \ref{app:quad_dyn}). Comparisons (see App. \ref{app:probs} for details) include na\"ive, uncalibrated MPC on the learned dynamics $\hat f$ without robust constraint tightenings, denoted as vanilla MPC (V-MPC), and a fixed hyper-ball variant of our method (CP-Ball), where $\covariance(z_k,v_k)$ in \eqref{eq:non_conform} is replaced with $I_{n_x}$, highlighting the benefit of state-dependent CP-Ellipsoid (our method) uncertainty sets. CP-Ball is an adaptation of \citet{DBLP:conf/l4dc/CheeSHP24}, using the L2-norm of errors as the non-conformity score. 
% Direct use of \citet{DBLP:conf/l4dc/CheeSHP24} was not feasible, since it requires larger horizon lengths for successful application, as it relies solely on online data by collecting the first 150 online data-points prior to deploying their conformal prediction algorithm. 
Direct use of \citet{DBLP:conf/l4dc/CheeSHP24} was not feasible, as it only collects online data points, requiring horizon lengths longer than 150 points for successful deployment. 
Training details, model architectures, and cost functions are in App.~\ref{app:training}. Our implementation uses Acados \citep{verschueren2020acadosmodularopensourceframework} and L4CasADi \citep{salzmann2024l4casadi} and is run on an M4 Pro MacBook (12 cores, 24 GB RAM). We report average MPC step time, prediction error $|f(x,u) - \hat f(x,u)|$, and minimum obstacle distance. All trials use horizon $T=15$ with $\alpha_k = 0.1/15$.

% In this section, we evaluate our approach on robust planning with learned 4D Dubins' and 12D quadcopter models. 
% To evaluate our method, we compare it against vanilla MPC (V-MPC) without robust constraint tightenings, quantifying the performance of na\"ive planning with learned dynamics. We also compare to a variant of our method, dubbed the fixed hyper-ball approach (CP-Ball), where the covariance matrix $\bm \Sigma(z_k,v_k)$ in \eqref{eq:non_conform} is replaced with the identity $I_{n_x}$. This allows us to assess the performance gain of using state-dependent hyper-ellipsoid (CP-Ellipsoid) uncertainty sets to inform SLS. For our experiments, we learn the dynamics of a 4D Dubins' car with 2 control inputs and a 12D quadcopter with 4 control inputs (see App.~\ref{app:car_dyn} and \ref{app:quad_dyn}, respectively).
% Training details, model architectures, and cost functions are provided in App.~\ref{app:training}. Our implementation uses PyTorch for learned models, Acados for optimization, and L4CasAdi for integration between Acados and CasADi. Experiments were run on an M4 Pro MacBook with 12 CPU cores and 24 GB of memory. Reported metrics include the average computation time per MPC step, average prediction error $\|f(x,u) - \hat f(x,u)\|$, and minimum obstacle distance. All experiments use a horizon length $T=15$ with $\alpha_k = 0.1/15$ at each timestep.


% To compare our method, we used vanilla-MPC without robust constraint tightenings (or SLS) to quantify the performance of naive planning with learned dynamics. We also compare this approach with using the identity matrix (replacing $\bm \Sigma(z_k,v_k)$ with $I_{n_x}$ in \eqref{eq:non_conform}) to produce fixed hyper-balls (fixed ball) to identify the performance improvement by using state-dependent hyper-ellipsoid (adaptive ellipsoid) uncertainty sets to inform SLS. 
% For our dynamic models, we learn the dynamics of a 4D Dubins' car with 2 control inputs and a 12D Quadcopter with 4 control inputs as defined in App. \ref{app:car_dyn} and \ref{app:quad_dyn}, respectively. To test our method on models which are not guaranteed to generalize well outside of the training domain we augment the dynamics in the out of distribution (OOD) regions. Lastly, we discuss training details and model architectures, and cost functions in App. \ref{app:training}. 

% We implement our algorithm using PyTorch for learned models, Acados for Optimization, and L4CasAdi for integration between Acados and L4CasAdi. To run our experiments, we used an M4 Pro MacBook Pro with 12 CPU cores and 24 GB of memory. Using our implementation, we report the average computation time for each MPC step, the average prediction error ($||f(x,u) -\hat{f}(x,u)||$, and the minimum obstacle distance. All our experiments were conducted with a horizon length $T=15$ and with $\alpha_k = 0.1/15$ for each timestep.

% \begin{itemize}
%     \item Fixed conformal prediction bound (we expect to have more violations on average, as the coverage guarantees only hold on average over a larger domain)
%     \item Naive planning with the learned dynamics without accounting for model error (just throw into Acados, without tightenings)
% \end{itemize}

% \subsection{Car}
\begin{figure}[htbp]
    \centering
    \vspace{-8pt}
    \includegraphics[width=\textwidth]{figs/ellip_tube.pdf}
    \vspace{-25pt}
    \caption{\textbf{OOD Car.} (Top) Minimum distance between prediction error and the ellipsoid edge is plotted for varying $\calib$ with a fixed start and goal; negative values (blue) indicate points inside the ellipsoid, i.e., valid coverage. (Bottom) Corresponding one-step $\theta$ tubes are shown. Smaller $\calib$ keeps prediction errors inside the ellipsoid more often and results in larger tubes.%(Top) We plot the minimum distance between the prediction error and the edge of the ellipsoid while varying the $\calib$ for a fixed start and goal. The distance is negative for points insside the ellipsoid and positive outside. (Bottom) We provide the corresponding one-step tubes for $\theta$. With smaller $\calib$ we find the prediction error remains within the ellipsoid more often and has larger tubes.
    }\vspace{-15pt}
    \label{fig:OOD_tubes}
\end{figure}
\paragraph{Dubins': In Domain (ID).} First, we validate that our approach maintains performance compared to baselines in ID scenarios without a major loss of computational efficiency. We used the car dynamics in \eqref{eq:car_dynamics} and selected start and goal states training distribution (equivalently the feasible set for the ID results). We define $\mathcal{F} = \{(x,u) \mid p_x \in [0,5], p_y \in [-5,5], \Vert [p_x, p_y]^\top -[2.5, 0]\Vert_2 \ge 1\}$, 

\begin{wrapfigure}{r}{0.35\textwidth}
    \centering
    \vspace{-5pt}
    \includegraphics[width=\linewidth]{figs/run_1_friction_car_truncated.pdf}
    \vspace{-0.85in}
    \caption{\textbf{Friction Car.} We plot all approaches and forecasted MPC steps \& tubes for CP-Ellipsoid.}
    \label{fig:friction_car}
    \vspace{-10pt}
\end{wrapfigure}
\vspace{-5pt}

\noindent with bounds on $p_x$ and $p_y$ and a 1-m radius obstacle centered at $(2.5, 0)$. Across 10 runs, all three approaches (V-MPC, CP-Ball, and CP-Ellipsoid) safely reached the goal while avoiding the obstacle; however, V-MPC does not perform constraint tightening and lacks safety guarantees. In Table \ref{tab:combined_results_modified}, CP-Ball and CP-Ellipsoid have larger obstacle proximity--a byproduct of the constraint tightening (Table \ref{tab:combined_results_modified}).
Finally, the one-step prediction error for $\hat f$ remains inside the ball and ellipsoid over $99.33\% = 100(1-\frac{0.1}{15})\%$ during execution--validating the coverage guarantees (Fig. \ref{fig:ellip_in_domain} in App. \ref{app:additional_results}).
% V-MPC is faster than CP-Ball and CP-Ellipsoid, since it does not spend time on reachability and constraint tightening. 
In the training domain, we expect V-MPC to have similar success rates as our method since $\hat f$ is accurate; this changes in OOD settings. We find that our approach avoids the obstacle ID, without a major slowdown from constraint tightening computations.
% Overall, this experiment suggests that our approach maintains performance ID without a major slowdown.

\vspace{-5pt}
\paragraph{Dubins': OOD.} To evaluate OOD performance, we train a dynamics and uncertainty model where  $p_y \in [-12, -6] \cup [6,12]$ and modify the dynamics to steer the car’s angle toward the obstacle for $p_y \in [-6,6]$ (see App. \ref{app:car_dyn}), while placing the feasible set of the rollout (same as the ID experiment) OOD. In this setting, we vary the calibration set size from 2,250 to 4,250 points (Fig. \ref{fig:OOD_tubes}) for CP-Ellipsoid, ensuring the start and goal are outside the training domain. 
\begin{wrapfigure}{r}{0.38\textwidth}
\vspace{-16pt}
    \centering
    \includegraphics[width=\linewidth]{figs/log_volume_histogram.pdf}
    \vspace{-25pt}
    \caption{\textbf{Active Uncertainty.} We plot the histogram of summed log-volumes for each forecasted MPC plan. With active uncertainty reduction, the distribution shifts left, indicating less uncertainty at runtime.%We plot the histogram of the sum of log-volumes for each forecasted MPC plan. Using the active uncertainty reduction approach we find the distribution of log-volumes shifts left, indicating this approach reduces encountered uncertainty during rollouts.
    }
\vspace{-10pt}
    \label{fig:friction_car}
\end{wrapfigure}
% \vspace{-10pt}
As shown in Fig.~\ref{fig:OOD_tubes}, the one-step prediction errors initially lie outside the predicted ellipsoid but gradually return within it as online calibration points are collected. Moreover, with fewer initial calibration points, the online adaptation occurs more rapidly, resulting in a higher coverage of errors in the ellipsoid. The plot also shows that the one-step tubes in $\theta$ are wider for smaller $\calib$, indicating that data scarcity increases conservativeness. The goal is reliably reached, suggesting that the SLS tubes, when informed by CP, enable robust OOD performance.
% Next, to evaluate fully OOD performance, we learn a dynamics and uncertainty model such that $p_y \in [-12, -6] \cup [6,12]$ and introduce dynamics that gradually steer the the car's angle towards the obstacle for $p_y \in [-6,6]$ (discussed in App. \ref{app:car_dyn}). We also ensure our start and end goals remain outside of the training domain, while maintaining the obstacle from our in domain experiment. In this experiment we changed size of the calibration set from $2,250$ points to $4,250$ points.
% In Figure \ref{fig:OOD_tubes}, we find that the one-step prediction errors start outside of the predicted ellipsoid and over-time come back into the ellipsoid due to the online collection of calibration points. Additionally, with smaller initial calibration points, we find that the online adaptation is quicker leading to greater percentage errors contained within the ellipsoid. The plot also demonstrates that the one-step $\theta$ tube remains wider for smaller $\calib$. Since the car was successful in reaching the goal this suggests that tube over-approximation in SLS coupled with conformal prediction enables robust performance out of domain.

% \begin{wrapfigure}{r}{0.4\textwidth}
% \vspace{-20pt}
%     \centering
%     \includegraphics[width=\linewidth]{figs/log_volume_histogram.pdf}
%     \vspace{-25pt}
%     \caption{\textbf{Active Uncertainty.} We plot the histogram of summed log-volumes for each forecasted MPC plan. With active uncertainty reduction, the distribution shifts left, indicating lower encountered uncertainty during rollouts.%We plot the histogram of the sum of log-volumes for each forecasted MPC plan. Using the active uncertainty reduction approach we find the distribution of log-volumes shifts left, indicating this approach reduces encountered uncertainty during rollouts.
%     }
% \vspace{-15pt}
%     \label{fig:friction_car}
% \end{wrapfigure}
\vspace{-5pt}
\paragraph{Dubins': Disjoint Training Domains (Friction Car).} 
To evaluate OOD obstacle avoidance under limited online data, we construct disjoint training regions (similar to the OOD setup) where  $p_y \in [-5, -2] \cup [2,5]$. The start and goal positions are randomly sampled from different regions, ensuring that all feasible trajectories must traverse OOD areas to reach the goal. In the OOD region, we apply the same steering modification as in the previous experiment, while adding $p_y$ dependent acceleration dynamics (see App.~\ref{app:car_dyn} for details).
% To assess OOD obstacle avoidance, with a limited number of online samples we constructed an experiment with disjoint training regions (similar to OOD experiment) with $p_y \in [-5, -2] \cup [2,5]$, while ensuring our randomly sampled start and end goals each start in different sets such that all paths would have to go out of distribution to successfully reach the goal. In the OOD region we enable attractive steering like the previous experiment and added friction through the training domain, and slippage factors OOD (as described in App. \ref{app:car_dyn}).
Across 10 runs, V-MPC failed twice, the CP-Ball failed once, and CP-Ellipsoid was consistently successful. These differences arise from reduced obstacle clearance (Table~\ref{tab:combined_results_modified}) observed in the CP-Ball and V-MPC cases. Fig.~\ref{fig:friction_car} shows that both CP-Ball and CP-Ellipsoid reach the goal, but CP-Ellipsoid maintains a greater distance from the obstacle, as evident from its forecasted trajectories and uncertainty tubes. Fig. \ref{fig:friction_car} also illustrates that the state-dependent ellipsoids expand along the direction of motion, as increased forecasted $y$-motion leads to higher $y$-direction uncertainty. A result from an extra run is in App.~\ref{app:additional_results}.
% Across 10 runs, the vanilla-MPC approach failed twice, the fix ball failed once, and the adaptive ellipsoid was always successful. This is due to decreased obstacle proximity (Table \ref{tab:mean_min_dist}) for the fixed ball and vanilla-MPC approaches. In Figure \ref{fig:friction_car}, we observe that both tube-based approaches successfully reach the goal, while the adaptive ellipsoid selects steps to maintain a greater distance to the obstacle (as seen by its forecasted steps and tubes). In this figure, we also observe that the state-dependent ellipsoids cause the tubes to grow in the direction of motion as the forecasted y-direction motion increases the y-direction uncertainty. We also provide a similar figure from another run in App. \ref{app:additional_results}.

% \vspace{-8pt}
\paragraph{Dubins': Active Uncertainty Reduction.}
Finally, to evaluate active uncertainty reduction, we trained $\hat{f}$ and $\covariance$ without sampling data in a circular region centered at $(p_x, p_y) = (2.5, 0.0)$ with a radius of 1 meter, and introduced friction and attractive steering dynamics within this unseen region. Thus the feasible set contains an OOD region, where rollouts are permitted to traverse. Despite this, the active uncertainty reduction method achieves lower prediction error by planning a longer path that remains in the training domain (Table~\ref{tab:active_uncert} App. \ref{app:additional_results}). Moreover, Fig.~\ref{fig:friction_car} shows the log-volume of the uncertainty tubes, showing that the active uncertainty cost (\eqref{eq:active_cost} in App.~\ref{app:uncert_reduc}) yields smaller tubes (Fig. \ref{fig:friction_car}, green) compared to the baseline without this cost (Fig. \ref{fig:friction_car}, orange).
% \paragraph{Active Uncertainty Reduction.} Lastly, to test active uncertainty reduction, we conduct an experiment where we train models without sampling points around a circular region centered at (2.5, 0.0) with a radius of 1 meter. We then introduce friction dynamics and attractive steering in this region. Since there is no obstacle between the start and goal state, the trajectory optimizer can plan through this region. As a result, the active uncertainty reduction approach has a lower prediction error by remaining in the training domain (Table \ref{tab:active_uncert}). Additionally, Figure \ref{fig:friction_car} provides the log volume of tubes and demonstrates a reduction in uncertainty by using the active uncertainty cost (described in App. \ref{app:uncert_reduc}).



% \begin{itemize}
%     \item Planning within domain (for random start/goals, just sample across the full training domain). Also include a version of this where we don't use the uncertainty search, so that the planning is much faster, to show that the reason why we are slow is because of a lack.
%     \item Show a version where we just plan for the full horizon (not using MPC), to argue that we can make it to the goal without replanning (sidestepping the need for real-time solve times).
%     \item Planning between two disjoint data regions (could either implement random start/goals, starts sampled from one region and goals sampled from the other; or just both could be sampled from the same region; just choose one and stick with it). Show a scatter plot which shows tube sizes increases in the uncertain domain. Show a heatmap of how the delta allocations change over time (that they increase with high velocity and dynamics uncertainty).
%     \item C-shaped domain, showing that our method learns to go to high-data-density regions, even if they lead to higher path cost (for random start/goals, just sample across the full training domain)
% \end{itemize}

% Add these to your LaTeX preamble if you don't have them
% \usepackage{booktabs}
% \usepackage{multirow}

% Add these to your LaTeX preamble if you don't have them
% \usepackage{booktabs}
\vspace{-4pt}
\begin{table}[h]
    \centering
  { % Start a local group to keep \tabcolsep change local
  \footnotesize  
  \centering
  % Modified Caption to reflect changes
  \caption{\looseness-1Mean and standard deviation of computation times and prediction errors and the mean minimum obstacle distance for vanilla-MPC (V-MPC), the CP-Ball (Ball), and CP-Ellipsoid (Ellipsoid). }\vspace{-10pt}
  \label{tab:combined_results_modified}
  \setlength{\tabcolsep}{2pt} % <-- Kept reduced spacing
  \begin{tabular}{@{}lccccccccc@{}}
    \toprule
    % Main metric headers, each spanning 3 sub-columns
    & \multicolumn{3}{c}{\textbf{Computation Time (ms)}} 
    % Updated L2 error header to show scaling
    & \multicolumn{3}{c}{\textbf{Prediction Error ($L_2$) ($\times 10^{-2}$)}} 
    & \multicolumn{3}{c}{\textbf{Min. Obstacle Dist. (m)}} \\
    % \cmidrule adds partial lines under the multi-column headers
    \cmidrule(lr){2-4} \cmidrule(lr){5-7} \cmidrule(lr){8-10}
    % Sub-headers for the methods
    \textbf{Model} 
    & \textbf{V-MPC} & \textbf{Ball} & \textbf{Ellipsoid} 
    & \textbf{V-MPC} & \textbf{Ball} & \textbf{Ellipsoid} 
    & \textbf{V-MPC} & \textbf{Ball} & \textbf{Ellipsoid} \\
    \midrule
    
    % --- Data rows ---
    
    Car 
    % Comp Time: Rounded to 3 decimal places
    & 53.3 $\pm$ 2.1 & 126.6 $\pm$ 8.4 & 126.7 $\pm$ 8.3 
    % L2 Error: Scaled by 100
    & 0.28 $\pm$ 0.54 & 0.26 $\pm$ 0.14 & 0.26 $\pm$ 0.14 
    % Min Dist: Std. Dev. removed
    & 1.152 & 1.202 & 1.207 \\
    
    Friction Car 
    % Comp Time: Rounded to 3 decimal places
    & 53.7 $\pm$ 2.6 & 131.1 $\pm$ 1.4 & 129.7 $\pm$ 9.0 
    % L2 Error: Scaled by 100
    & 1.94 $\pm$ 2.87 & 1.92 $\pm$ 2.80 & 1.95 $\pm$ 2.69 
    % Min Dist: Std. Dev. removed
    & 1.315 & 1.376 & 1.464 \\
    
    Quadcopter 
    % Comp Time: Rounded to 3 decimal places
    & 63.1 $\pm$ 4.6 & nan $\pm$ nan & 222 $\pm$ 20.9 
    % L2 Error: Scaled by 100
    & 3.89 $\pm$ 5.49 & nan $\pm$ nan & 3.30 $\pm$ 4.17 
    % Min Dist: Std. Dev. removed
    & 0.804 & nan & 0.809 \\
    \bottomrule
  \end{tabular}\vspace{-5pt}
  } % End the local group
\end{table}



\paragraph{12D Quadcopter}
\begin{SCfigure}
\vspace{-10pt}
    \centering
    \includegraphics[width=0.65\textwidth]{figs/quad_trajectory_light.pdf}
    \vspace{-7pt}
    \caption{\textbf{Quadcopter.} A rollout of the quadcopter with V-MPC and CP-Ellipsoid. CP-Ellipsoid maintains sufficient proximity to avoid crashing into the obstacle, while V-MPC crashes.}
    \label{fig:quad_trajectory}
\end{SCfigure}
\vspace{-10pt}
To evaluate the scalability of our method to higher-dimensional systems, we train a 12D quadcopter dynamics model. Using this model, we plan trajectories such that the spatial coordinates $(p_x, p_y, p_z)$ remain outside a sphere centered at $(2.5, 2.5, 2.5)$ with a radius of 0.8 meters -- leaving a 0.2-meter-thick shell around the sphere that is out of distribution and characterized by increased downward acceleration. To improve the performance of the uncertainty model, we modify its weighting by blending the covariance matrix with the identity matrix to reduce conservativeness, i.e., using $((1-\tau) I + \tau \covariance(z,v))$ instead of $\covariance(z,v)$ in \eqref{eq:ellip_prediction_set}. With this approach, the CP-Ellipsoid method successfully planned 8 out of 10 trajectories (one failure was due to numerical instability, which did not result in collision), whereas V-MPC approach succeeded in only 4 out of 10 runs. While this demonstrates improved robustness over V-MPC, the failed case suggests faster tube expansion may be necessary to ensure consistent collision avoidance. Fig.~\ref{fig:quad_trajectory} shows an example quadcopter trajectory, demonstrating that our CP-informed tubes maintain safe obstacle proximity.


% To demonstrate the effectiveness of our framework on larger models, we train a 12D quadcopter dynamics model. Using this model, we planned trajectories such that the spatial coordinates $p_x, p_y, \text{ and } p_z$ were outside of a ball centered at (2.5, 2.5, 2.5) and a radius of 0.8 meters - leaving a 0.2 meter shell around this ball that was out-of-domain with a greater falling acceleration. 
% To enable our uncertainty model to perform better, we weighted it with the covariance matrix with an identity matrix to reduce conservativeness ($(1-\tau) I_{12} + \tau \bm \Sigma(z,v)$ instead of $\bm\Sigma(z,v)$ in \eqref{eq:ellip_prediction_set}). Using this approach, the quadcopter successfully planned $8/10$ trajectories (with one trajectory failing due to numerical instability, but did not hit the obstacle) with the adaptive ellipsoid and $4/10$ trajectories with the vanilla MPC approach. While we improve over vanilla MPC, the failed run suggests that tubes should grow faster to ensure successful collision avoidance. Lastly, Figure \ref{fig:quad_trajectory} shows a successful quadcopter trajectory, demonstrating how the tubes--informed by conformal prediction--ensure proximity is maintained from the ellipsoid. 
% Planning inside a thin tube around one nominal trajectory. Collect various trajectory optimization solutions around the one nominal (by randomly perturbing start, goal, or by perturbing from intermediate points along the trajectory). Also perturb the velocity, angle states, to make sure we have enough coverage.

% Another way to collect data:
% \begin{itemize}
%     \item Compute one reference trajectory using trajectory optimization between start and goal, that avoids an obstacle.
%     \item Do one-step random sampling around it: randomly sample a state on the reference; sample a Gaussian/uniform perturbation around that state, a random perturbation of the control input at that point. Use this to augment the dynamics training dataset.
% \end{itemize}
\vspace{-11pt}
\section{Conclusion}
\vspace{-4pt}

In this work, we developed an efficient framework for robust out-of-distribution MPC with learned dynamics, using WCP and SLS. We provided theoretical guarantees ensuring that the rollout satisfy constraints with high probability, and introduced an uncertainty reduction cost function to encourage the system to remain within the training domain when possible. Through experiments on variants of the 4D Dubins car and 12D quadcopter, we demonstrated that the \methodnospace{} framework achieves robust performance even when using learned dynamics that generalize poorly OOD.
% In this work we, developed a fast framework for robustly planning out of distribution with learned dynamics via WCP and SLS. We also provided theoretical guarantees of a rollout remaining in a feasible a set. To improve performance, we also provided an uncertainty reduction cost function to prioritize remaining in the training domain when possible. Through our experiments on variants of the 4D Dubins' car and 12 Quadcopter we demonstrated the \methodnospace's framework ability to perform robustly despite using learned dynamics models which do not generalize well outside of the training domain.  
\acks{We thank Devesh Nath for insightful discussions and feedback on this work.}

\bibliography{l4dc2026-sample}
\appendix

\newpage
\section{Proofs:}\label{app:proofs}
\subsection{Proof of Theorem \ref{thm:cov_gap_thm} and Corollary \ref{cor:cov_gap_infty}}
\begin{proof} \textbf{(Theorem \ref{thm:cov_gap_thm})}
From Theorem 2 and Lemma 1 in \citet{barber2023conformal}, the coverage gap for independent data is bounded by:
\begin{equation}\label{eq:init_barber}
\text{Coverage Gap} \le \sum_{i=1}^{N_\textrm{calib}} \tilde{w}_i \cdot d_{TV}(S^{i,k}, S^{k,k}) \le \sum_{i=1}^{N_\textrm{calib}} \tilde{w}_i \cdot 2 d_{TV}(S_{i,k}, S_{k,k})
\end{equation}
Given the normalized weights are $\tilde{w}_i = \frac{w_i}{1 + \sum_{j=1}^{N_\textrm{calib}} w_j}$, and $w_i = \rho^{d_i}$ for $d_i = ||[z_k\T,v_k\T] - [x_i\T,u_i\T]||_2$
We can substitute into Equation \ref{eq:init_barber}, and prove the \textit{Tight Bound} in Theorem \ref{thm:cov_gap_thm}:
\begin{equation}\label{eq:sub_barber}
\text{Coverage Gap} \le \sum_{i=1}^{N_\textrm{calib}} \left( \frac{\rho^{d_i}}{1 + \sum_{j=1}^{N_\textrm{calib}} \rho^{d_j}} \right) \cdot 2\epsilon \cdot d_i
\end{equation}
Without loss of generality, we can rearrange the sum to create a decaying series $\{\rho^{d_i}\}_{i=1}^{N_\textrm{calib}}$, where $d_i$ is increasing. While this series is not purely geometric, we can bound it by two geometric series. Let $d_\text{min} = \min\{1 \leq i \leq N_\textrm{calib}-1 | d_{i+1} - d_i\}$ and $d_\text{max} = \max\{1 \leq i \leq N_\textrm{calib}-1 | d_{i+1} - d_i\}$. Then, $\rho^{d_1}\rho^{d_\text{max}(i-1)} \leq \rho^{d_i}  \leq \rho^{d_1}\rho^{d_\text{min}(i-1)}$. Thus, we can bound Equation \ref{eq:sub_barber}, 
$$
\text{Coverage Gap} \le \sum_{i=1}^{N_\textrm{calib}} \left( \frac{\rho^{d_i}}{1 + \sum_{j=1}^{N_\textrm{calib}} \rho^{d_j}} \right) \cdot 2\epsilon \cdot d_i \leq 2\epsilon \left( \frac{\sum_{i=1}^{N_\textrm{calib}}\rho^{d_1}p^{d_\text{min}(i-1)}\cdot (d_1 + d_\text{max}(i-1))}{1 + \sum_{j=1}^{N_\textrm{calib}} \rho^{d_1}\rho^{d_\text{max}(j-1)}} \right) 
$$
Then observe the sum of the numerator, 
\begin{align}
N_\text{sum} &= \sum_{i=1}^{N_\textrm{calib}}\rho^{d_1}p^{d_\text{min}(i-1)}\cdot (d_1 + d_\text{max}(i-1)) \\
&=\rho^{d_1} \left[ d_1 \left( \frac{1 - \rho^{N_\textrm{calib} \cdot d_\text{min}}}{1 - \rho^{d_\text{min}}} \right) + d_\text{max} \left( \frac{\rho^{d_\text{min}} - N_\textrm{calib} \rho^{n \cdot d_\text{min}} + (N_\textrm{calib}-1)\rho^{(N_\textrm{calib}+1)d_\text{min}}}{(1-\rho^{d_\text{min}})^2} \right) \right]
\end{align}

Similarly, the sum the of the denominator,
\begin{align}
    D_\text{sum} &= 1 + \sum_{j=1}^{N_\textrm{calib}} \rho^{d_1}\rho^{d_\text{max}(j-1)}\\
    &= 1 + \rho^{d_1} \left( \frac{1 - \rho^{N_\textrm{calib} \cdot d_\text{max}}}{1 - \rho^{d_\text{max}}} \right)
\end{align}

Then, the coverage gap is bounded by,
\begin{small}
\begin{align}
    &\text{Coverage Gap} \le 2\epsilon\frac{N_\text{sum}}{D_\text{sum}} \\
    &\le 2 \epsilon \left[{d_1 \left( \frac{1 - \rho^{N_\textrm{calib} \cdot d_\text{min}}}{1 - \rho^{d_\text{min}}} \right) +
    d_\text{max} \left( \frac{\rho^{d_\text{min}} \overbrace{- N_\textrm{calib} \rho^{N_\textrm{calib} \cdot d_\text{min}} + (N_\textrm{calib}-1)\rho^{(N_\textrm{calib}+1)d_\text{min}}}^{\leq 0}}{(1-\rho^{d_\text{min}})^2} \right)}\right]\underbrace{\frac{1 - \rho^{d_\text{max}}}{1 - \rho^{N_\textrm{calib} \cdot d_\text{max}}}}_{\leq 1}\label{eq:n_sum_d_sum}\\
    &\le 2 \epsilon \left[  \frac{d_1 }{1 - \rho^{d_\text{min}}} + \frac{d_\text{max}\rho^{d_\text{min}}}{(1-\rho^{d_\text{min}})^2}\right].\label{eq:final_bound} 
\end{align}
Hence, we have proved the \textit{Interpretable Bound}.
\end{small}
\end{proof}


For large $N_\textrm{calib}$, (i.e., as $N_\textrm{calib}\to\infty$), we can achieve a tighter upper bound
\begin{corollary}\label{cor:cov_gap_infty} Under the same assumptions of Theorem \ref{thm:cov_gap_thm} we have that as $N_\textrm{calib}\to \infty$ the coverage gap is more tightly bounded by,
\begin{equation}\text{Coverage Gap}\le 2\epsilon
\left[ \frac{d_1 }{1 - \rho^{d_\text{min}}} + \frac{d_\text{max}\rho^{d_\text{min}}}{(1-\rho^{d_\text{min}})^2}\right](1-\rho^{d_\text{max}})
\end{equation}
\end{corollary}

\begin{proof}\textbf{(Corollary \ref{cor:cov_gap_infty})}
Taking Equation \ref{eq:n_sum_d_sum} as $n\to\infty$ we complete the proof. 
\end{proof}

\subsection{Proof of Theorem \ref{thm:tube_gap}}
\begin{proof}\textbf{(Theorem \ref{thm:tube_gap})}
    Using Equation \ref{eq:init_barber}, we have that 
    \begin{equation}\label{eq:init_barber_2}
    \text{Coverage Gap} \le \sum_{i=1}^{N_\textrm{calib}} \tilde{w}_i \cdot d_{TV}(\tilde{S}^{i,k}, \tilde{S}^{k,k}) \le \sum_{i=1}^{N_\textrm{calib}} \tilde{w}_i \cdot 2 d_{TV}(S_{i,k}, \tilde{S}_{k,k}),
    \end{equation}
    where $(s_{1,k}, s_{2,k}, s_{i-1,k}, \tilde{s}_{k,k}, s_{i+1,k}, \ldots, s_{N_\textrm{calib},k}) \sim \tilde{S}^{i,k}$, $(s_{1,k}, \ldots, s_{N_\textrm{calib},k}, \tilde{s}_{k,k}) \sim \tilde{S}^{k,k}$, $\tilde{s}_{k,k} = s_k(x_k,u_k,f(x_k,u_k))$ is unknown due to uncertainty in $f(x_k,u_k)$, and $\tilde{s}_{k,k} \sim \tilde{S}_{k,k}$ is the distribution of the non-conformity score corresponding to the observed state and control $(x_k,u_k)$. Then, using the triangle inequality we have that,
    \begin{equation}\label{eq:split_into_sums}
    \text{Coverage Gap} \le \sum_{i=1}^{N_\textrm{calib}} \tilde{w}_i \cdot d_{TV}(\tilde{S}^{i,k}, \tilde{S}^{k,k}) \le \underbrace{\sum_{i=1}^{N_\textrm{calib}} \tilde{w}_i \cdot 2 d_{TV}(S_{i,k}, S_{k,k})}_{(*), \text{Bounded in Eq. \ref{eq:final_bound}}} + \underbrace{\sum_{i=1}^{N_\textrm{calib}} \tilde{w}_i \cdot 2 d_{TV}(S_{k,k}, \tilde{S}_{k,k})}_{(**)}.
    \end{equation}
    Observing that the left sum $(*)$ is already bounded by Theorem \ref{thm:cov_gap_thm}, we proceed to bound the right sum $(**)$. Then, realizing the sum of normalized weights is $\le 1$, we obtain, 
    \begin{equation}\label{eq:near_by_cov}
        \sum_{i=1}^{N_\textrm{calib}} \tilde{w}_i \cdot 2 d_{TV}(S_{k,k}, \tilde{S}_{k,k}) = 2 d_{TV}(S_{k,k}, \tilde{S}_{k,k})\sum_{i=1}^{N_\textrm{calib}} \tilde{w}_i \le 2 d_{TV}(S_{k,k}, \tilde{S}_{k,k})\
    \end{equation}
     Then, using the Lipschitz-type bound, 
    \begin{equation}
        2 d_{TV}(S_{k,k}, \tilde{S}_{k,k}) \le 2\hat{\epsilon}||[x_k\T,u_k\T]-[z_k\T,v_k\T]||_2 \le 2\hat{\epsilon}M(\mathcal{R}_k) =: \gamma({\mathcal{R}_k}),
    \end{equation} 
    where, $(x_k,u_k) \in \mathcal{R}_k$, $M(\mathcal{R}_k)$ is the maximum axis length of the tube, and $\hat{\epsilon} \le \epsilon$ is the local Lipschitz constant in the tube. The remaining proof follows substituting the bound, $\gamma({\mathcal{R}}_k)$ for $(**)$ in \eqref{eq:split_into_sums}.
\end{proof}

\subsection{Proof Sketch of Theorem \ref{thm:total_coverage} and Theorem \ref{cor:rollout_prob}}
\begin{proofsketch} \textbf{(Theorem \ref{thm:total_coverage})}
Using Theorem \ref{thm:tube_gap} we get the probability of error coverage for each step in the rollout. Using the approach of Proposition 2 and Theorem 1 from \citet{DBLP:conf/l4dc/CheeSHP24} and \citet{DBLP:journals/ral/LindemannCSP23} respectively we complete the proof. 
\end{proofsketch}
\begin{proofsketch}\textbf{(Theorem \ref{cor:rollout_prob})}
\textbf{(a)} Since SLS provides guarantees conditioned on valid disturbance bounds, Theorem \ref{thm:total_coverage} directly implies the result in Corollary \ref{cor:rollout_prob}(a).

\textbf{(b)} Instead of applying the union bound approach across the forecasted MPC plan, we bound the intersection probability of the first-step errors remaining in the predicted ellipsoids for all $R$ planning steps. Then, using the same approach as Theorem \ref{thm:tube_gap} we get: 
\begin{equation}
        \Pr[\bigcap_{k=1}^{T}~\varepsilon(z_1^{(q)}, v_1^{(q)}) \in \mathcal{E}(z_1^{(q)}, v_1^{(q)})] \geq 1- \sum_{q=1}^{R}\sigma_1^q,
\end{equation}
Notably, we omit the tubes since the tubes at the first step of each replan, $\mathcal{R}_1^{q} = \{(z_1^{(q)}, v_1^{(q)})\}$ since the planner knows the initial state. Then, using the same approach as \textbf{(a)} we get the result in \textbf{(b)}.  
\end{proofsketch}


\section{Dynamics Models}
Below we discuss the dynamics models used in our experiments. We provide the continuous time dynamics which were used to train our models. We obtained the discrete model by using Euler Discretization. 
\subsection{Car Dynamics Models}\label{app:car_dyn}
\begin{equation}\label{eq:car_dynamics}
\begin{bmatrix}
    \dot{p}_x \\
    \dot{p}_y \\
    \dot{\theta} \\
    \dot{v}
\end{bmatrix}
=
\begin{bmatrix}
    v \cos(\theta) \\
    v \sin(\theta) \\
    0 \\
    0
\end{bmatrix}
+
\begin{bmatrix}
    0 & 0 \\
    0 & 0 \\
    1 & 0 \\
    0 & 1
\end{bmatrix}
\begin{bmatrix}
    \omega \\
    a
\end{bmatrix}
\end{equation}
Equation \ref{eq:car_dynamics} describes the control-affine dynamics models for Dubins' car, where $p_x$ and $p_y$ are the car's x and y positions. $\theta$ and $v$ are the direction and velocity of the car. Lastly, the control inputs directly steer ($\omega$) and accelerate ($a$) the car. The dynamics model in  \eqref{eq:car_dynamics} was used strictly for the in-distribution experiments. 

\paragraph{OOD Experiment Dynamics} For the out of distribution (OOD) Dubins' car experiment we modified the $\theta$ dynamics by adding attractive steering towards the obstacle centered at (2.5, 0) in the region without training data,
\begin{equation}\label{eq:theta_mod}
    \dot\theta = \omega_\text{attr}\1[|p_y| < 6] + \omega 
\end{equation}

\begin{subequations}\label{eq:w_attr}
\begin{align}
    d^2 &= (p_x - 2.5)^2 + p_y^2 \\
    \theta_{a} &= \text{atan2}(p_y,p_x) \\
    \Delta\theta_{norm} &= \text{atan2}(\sin(\theta_{a} - \theta), \cos(\theta_{a} - \theta)) \\
    \omega_{attr} &= \left( \frac{k_\text{attr}}{d^2} \right) \Delta\theta_{norm},
\end{align}
\end{subequations}
where the attractive steering dynamics are defined in Equation \ref{eq:w_attr}, and $\text{atan2}$ is the arctangent function accounting for the x-y quadrant, and $\1[\cdot]$ is the indicator function. For the OOD experiment we set $k_\text{attr}$, the strength of the attractive force to -0.5 (positive values repel). Notably, the strength of the force increases as the car gets closer to the obstacle. 

\paragraph{Active Uncertainty Experiment Dynamics} Similar to the OOD dynamics we include attractive steering (with $k_\text{attr} = -0.5$, we modify the the $v$ dynamics of the car such that, 
\begin{subequations}
    \begin{align}
    &\dot\theta = \tilde{\omega}_\text{attr}\1[d^2 < 1] + \omega \\
    &\dot v =  -(1 + \cos(\pi \sqrt{d^2}))\1[d^2 < 1]+ a \\
    &\tilde\omega_{attr} = \left( \frac{k_\text{attr}}{d^2 + 0.1} \right) \Delta\theta_{norm},
    \end{align}
\end{subequations}
where $d^2$ and $\Delta\theta_{norm}$ are defined in Equation \ref{eq:w_attr}. We modify the denominator of $\omega_\text{attr}$ to prevent a singularity at (2.5, 0.0)--this was not a problem in the OOD case since there was an obstacle present at (2.5, 0.0). The augmented velocity dynamics add a friction term to reduce the speed of the car in this region.

\paragraph{Friction Car Dynamics} Lastly, we modified Dubins' car such that the velocity dynamics were modified over the whole state-space, while including attractive steering (with $k_\text{attr}=-1.5$): 
\begin{subequations}
    \begin{align}
    &\dot\theta = {\omega}_\text{attr}\1[|p_y| < 2] + \omega \\
    &\dot v =  a-\text{sign}(v)\left(0.1\cos(\frac{2\pi}{5}p_y)) + 0.1 - a_\text{slip}\1[|p_y| < 2]\right)\\
    &a_\text{slip} = 0.7e^{-2p_y^2} 
    \end{align}
\end{subequations}
In this model, we include a slippage term to increase the velocity of the car for $-2 < p_y < 2$, while providing a changing friction term over the whole state-space. 

\subsection{12D Quadcopter Dynamics Model}\label{app:quad_dyn}
\begin{equation}
\begin{bmatrix}
    \dot{p}_x \\
    \dot{p}_y \\
    \dot{p}_z \\
    \dot{\psi} \\
    \dot{\theta} \\
    \dot{\phi} \\
    \ddot{p}_x \\
    \ddot{p}_y \\
    \ddot{p}_z \\
    \dot{p} \\
    \dot{q} \\
    \dot{r}
\end{bmatrix}
= f(x,u) :=
\begin{bmatrix}
    \dot{p}_x \\
    \dot{p}_y \\
    \dot{p}_z \\
    q \sin(\phi)/\cos(\theta) + r \cos(\phi)/\cos(\theta) \\
    q \cos(\phi) - r \sin(\phi) \\
    p + q \sin(\phi)\tan(\theta) + r \cos(\phi)\tan(\theta) \\
    \frac{u_1}{m} (\sin(\phi)\sin(\psi) + \cos(\phi)\cos(\psi)\sin(\theta)) \\
    \frac{u_1}{m} (\cos(\psi)\sin(\phi) - \cos(\phi)\sin(\psi)\sin(\theta)) \\
    g + \frac{u_1}{m} (\cos(\phi)\cos(\theta)) -a_\text{fall}\\
    \frac{I_y - I_z}{I_x} qr + \frac{u_2}{I_x} \\
    \frac{I_z - I_x}{I_y} pr + \frac{u_3}{I_y} \\
    \frac{I_x - I_y}{I_z} pq + \frac{u_4}{I_z}
\end{bmatrix}
\end{equation}
\begin{align}
    a_\text{fall} &= \left(0.1 + 0.1\cos(2\pi d^2 -\pi)\right)\1[{d^2 < 1}]\\
    d^2 &= (p_x-2.5)^2 + (p_y-2.5)^2 + (p_z-2.5)^2
\end{align}


For our experiments, we used the 12D Quadcopter dynamics model described in Equation \ref{app:quad_dyn} \citep{sabatino2015quadrotor}. $p_x, p_y, \text{ and } p_z$ describe the x, y, and z positions of the quadcopter while $\dot p_x, \dot p_y, \text{ and } \dot p_z$ describe the x, y, and z velocities. ${\psi}, {\theta}, \text{ and } {\phi}$ are the roll, pitch, and yaw. Similarly, $p, q, \text{ and } r$ are the roll rate, pitch rate, and yaw rates, respectively are directly controlled by inputs $u_2, u_3, \text{ and } u_4$. The control input, $u_1$, directly impacts the thrust--z-direction acceleration. Lastly, the constants $g$, $m$, $I_x$, $I_y$, and $I_z$ are the acceleration due to gravity, the mass of the quadcopter, and the principle moments of inertia. For our experiments, we let $g=-9.81 m/s^{2}$, $m = 1.0 kg$, $I_x = 0.5 kg\cdot m^2$, $I_y = 0.1 kg\cdot m^2$, and $I_z = 0.3 kg\cdot m^2$. Notably, the expression for $\ddot{p}_z$ includes $a_\text{fall}$ to introduce a falling term for OOD dynamics. The falling dynamics are not encountered while training the model as described in App. \ref{app:exp_data}. 

\section{Baseline Problem Definitions}\label{app:probs}

We define the baselines used in Sec. \ref{sec:results}. First, we define Vanilla MPC (V-MPC) \eqref{eq:vanilla_mpc}, which is a nominal na\"ive MPC planning under the learned dynamics without any constraint tightenings or model error calibration:
\begin{subequations}\small
\label{eq:vanilla_mpc}
\begin{align}
\min_{\substack{\mathbf{z}, \mathbf{v}}}\quad &  J(\mathbf{z},\mathbf{v}) + J_{\mathcal{X}_\textrm{f}}(\Z,\V)\\[-6pt]
\text { s.t. }\quad\quad &  {z}_{k+1}= \hat{f}(z_k,v_k),~ {z}_1 = \bar x_0,\quad\quad \forall k = 1,\ldots, T,\\
&g_{i}(z_k,v_k) + b_{i}  \leq 0,\\
&  \quad\quad \forall i =1, \ldots, n_c,\quad \forall k= 1,\ldots,T. \nonumber
\end{align}
\end{subequations}

We also define CP-Ball \eqref{eq:cp_ball}, which is a variant of our method which computes CP-calibrated model error bounds, except using a spherical, constant error bound representation, which is more conservative than our ellipsoidal bound:
\begin{subequations}\small
\label{eq:cp_ball}
\begin{align}
\min_{\substack{\mathbf{\Phi}_{\mathbf{x}},\mathbf{\Phi}_{\mathrm{u}},\mathbf{z}, \mathbf{v}}}\quad &  J(\mathbf{z},\mathbf{v})  + J_{\mathcal{X}_\textrm{f}}(\Z,\V) +  \tilde H_0(\Px, \Pu),\\[-6pt]
\text { s.t. }\quad\quad &  {z}_{k+1}= \hat{f}(z_k,v_k),~ {z}_1 = \bar x_0,\quad\quad \forall k = 1,\ldots, T,\\
&\Px_{k+1,j} =\bm{\nabla}_x \hat{f}\left(z_k, v_k\right)  \Px_{k,j} + \bm{\nabla}_u \hat{f}\left(z_k, v_k\right) \Pu_{k,j},\ \forall j = 1,\ldots, T, \forall k = j+1,\ldots,T,\\
&\Px_{j+1,j} = V,\quad \forall j = 1,\ldots, T,\\%~~\text{(Equation \ref{eq:ball_form})}
& \textstyle\sum_{j=1}^{k}\big\Vert[\nabla_xg_{i}(z_k,v_k)\Px_{k,j}, \nabla_ug_{i}(z_k,v_k)\Pu_{k,j}]\big\Vert_{2} + g_{i}(z_k,v_k) + b_{i} + g_i^\textrm{lin}(z_k,v_k)  \leq 0,\\
&  \quad\quad \forall i =1, \ldots, n_c,\quad \forall k= 1,\ldots,T. \nonumber
\end{align}
\end{subequations}

\noindent\textbf{Surrogate Cost:} All SLS formulations use \eqref{eq:regulizer} to minimize the volume of the tubes--a surrogate for uncertainty, where $\Px$ and $\Pu$ collects all $\Px_{k,j}$ and $\Pu_{k,j}$ respectively, and $P^{\frac{1}{2}}, Q^{\frac{1}{2}} \in \mathbb{R}^{n_x \times n_x}, R^{\frac{1}{2}}\in\mathbb{R}^{n_u \times n_u}$ are symmetric positive definite matrices.

\begin{equation}  \label{eq:regulizer}\small 
\tilde H_0(\Px,\Pu)\defeq 
\textstyle\sum_{j=0}^{T-1} \left( \|\Qf^{\frac{1}{2}} \Px_{T,j}\|_\F^2  + \sum_{k=j}^{T-1}  \left(\| Q^{\frac{1}{2}} \Px_{k,j}\|_\F^2+ \| R^{\frac{1}{2}} \Pu_{k,j}\|_\F^2 \right)\right),
\end{equation}


\section{Experiment Setup}\label{app:training}
\subsection{Data Sampling}\label{app:exp_data}
\paragraph{Dubins' Car \& Friction Car:} To train the in-distribution Dubins' car dynamics and uncertainty models, we applied uniform sampling to the following state-control space:
\begin{equation}\label{eq:in_dist_car}[p_x, p_y, \theta, v] \times[\omega, a]\in [0, 5] \times [-5, 5] \times [-\pi, \pi] \times [-10, 10] \times [-10, 10] \times [-10, 10]. \end{equation}
Similarly, for the out-of-distribution (OOD) car, we instead sampled $p_y \in [-12, -6] \cup [6,12]$ with the remaining state and controls remaining the same as \eqref{eq:in_dist_car}. For the friction car model, we sampled $p_y \in [-5, -2] \cup [2,5]$, while using \eqref{eq:in_dist_car} for the remaining dimensions. Lastly, for our active uncertainty experiments, we also sampled from Equation \ref{eq:in_dist_car}, and enforced that $(p_x - 2.5)^2 + (p_y)^2 > 1$ to create an OOD region. For both the OOD car and active uncertainty experiments, we ensured our learned model was restricted to training on areas where modifications did not impact the car dynamics model. 

Using the sampled points, we computed the ground truth using the expert dynamics model, as formulated in Equation \ref{eq:car_dynamics}, to create our dynamics training dataset (1,000,000 points), uncertainty training dataset (1,000,000 points), and calibration datasets (varying numbers of points).

\paragraph{12D Quadcopter:} To get informative data points to train our neural-network model, we collected expert trajectories such that the states satisfied,
\begin{subequations}\label{eq:opt_expert_quad}
\begin{align}
p_x, p_y, p_z &\in [0,5]\text{ s.t. } (p_x-2.5)^2 + (p_y-2.5) + (p_z-2.5)^2 > 1 \\
    \psi, \theta, \phi &\in [-\pi/4, \pi/4]\\
\dot{p}_x, \dot{p}_y, \dot{p}_z  &\in [-5, 5]\\
p,q,r &\in [-2, 2]\\
u_1 &\in [0, -2gm],~g = -9.81 m/{s^2}, m=1.0 kg \\
u_2, u_3, u_4 &\in [-0.1, 0.1]
\end{align}   
\end{subequations}
to create an OOD region centered at (2.5, 2.5, 2.5) with a radius of 1. Additionally, the goal position was $[p_x = 4.5, p_y = 2.5, p_z = 2.5]$, while the start positions was arbitrarily sampled from Equation \ref{eq:opt_expert_quad}. Our discrete timestep was $0.05s$ and our trajectory length was 100 timesteps. We sampled 100,000 trajectories and included state-control pairs from trajectories that the optimizer (IpOpt) was able to solve. This resulted in a dynamic training and uncertainty-trained set with $4,702,400$ points each, leaving up to $100,000$ points for the calibration dataset.

\subsection{Model Architecture and Training}
For our learned model (dynamics and uncertainty), we use an MLP with one hidden layer and the tanh activation function. Our uncertainty model was trained to output Cholesky factors, ensuring the matrices were positive-definite and facilitating faster performance in our optimization problem. 

For Dubins' car and the friction car dynamics models, we used $4,096$ and $2,048$ hidden nodes for the dynamics and uncertainty models, respectively.  
Similarly, for the 12D quadcopter, we used $1,024$ hidden nodes for both the dynamics and uncertainty models, and were trained for 200 epochs. Lastly, we used a learning rate of $10^{-4}$ and $10^{-5}$ for the dynamics and uncertainty models, respectively.

\subsection{Optimization Parameters and Cost Functions}\label{app:lqr_app}
For the SLS forward propagation and constraint tightenings, we used $\bar{P}^{\frac{1}{2}} = 10^6I_{n_x}$, $\bar{Q}^{\frac{1}{2}} = 10^6I_{n_x}$ and $\bar{R}^{\frac{1}{2}} = 10^6I_{n_u}$, as found in Equation \ref{eq:regulizer}. For trajectory optimization, we considered the following LQR cost, 
\begin{equation}\label{eq:lqr}
J(\Z,\V) :=  
      \sum_{k=1}^{T-1} \left( z_k\T Q  z_k + (z_{k+1} - z_k)\T Q_{s}  (z_{k+1} - z_k) + u_k\T R {u}_k \right),
\end{equation}
\begin{equation}\label{eq:lqr_terminal}
    J_{\mathcal{X}_\textrm{f}}(\Z,\V) := (z_T - z_\text{goal})\T Q_f (z_T - z_\text{goal})
\end{equation}
where $Q_f, Q_s, Q \in \mathbb{R}^{n_x \times n_x}$ are positive-semi-definite matrices, $R\in\mathbb{R}^{n_u \times n_u}$ is a positive-definite matrix, and $z_\text{goal}$ is the goal state. For all dynamics models, we used diagonal matrices, represented as $\text{diag}(a_1,a_2, \dots, a_n)$ and $\text{block\_diag}(\dots)$ representing block diagonal matrices. For Dubins' car (excluding the active uncertainty experiment) and the friction car, we let $Q_f = \text{diag}(1,1,0,1)$, $Q = Q_s = 0_{4 \times 4}$ (the $4 \times 4$ zero matrix), and $R = \text{diag}(0.1, 0.1)$. For the Quadcopter, $Q_f = \text{block\_diag}(I_3, 0_{3\times3}, I_3, 0.2I_3)$, $Q = \text{block\_diag}(0_{3\times3}, 0_{3\times3}, 0.01I_3,0_{3\times3})$, $Q_s = 0_{12\times12}$, and $R = \text{diag}(0.01, 0.1, 0.1, 0.1)$. 

Lastly, regarding conformal prediction parameters, we used $\rho = 0.97$ for the car model and $\rho = 0.8$ for the quadcopter model, respectively. Outside of the OOD car experiment, we used $10,000$ calibration points for the car models and $30,000$ calibration points for the quadcopter model. For all the models we set $\alpha_k = \frac{0.1}{15}$ s.t. the sum total of $\alpha$ would be 0.1.

\subsection{Active Uncertainty}\label{app:uncert_reduc} To enable active uncertainty reduction, we sought to minimize the distance to the positional-state variables (i.e., $p_x$ and $p_y$) from $\calib$. To achieve efficient performance, we used K-means clustering \cite{hastie2009elements} to find a representative set of $L$ points, $\tilde{\mathcal{D}} = \{({p_x}_i, {p_y}_i)\}_{i=1}^{L}$. Then, we minimized the following cost for Dubins' car,
\begin{align}
    &J_\text{active}(\Z,\V)\nonumber\\
    &=\sum_{k=0}^{T-1} \exp\left[-\frac{1}{L + G}\left(Ge^{-\beta||({p_x}_\text{goal}, {p_y}_\text{goal}, v_\text{goal}) - ({p_x}_k, {p_y}_k, v_k)||^2} + \sum_{j=1}^{L}e^{-\beta||({p_x}_j, {p_y}_j) - ({p_x}_k, {p_y}_k)||^2}\right)\right], \label{eq:active_cost}
\end{align}
where ${p_x}_\text{goal}, {p_y}_\text{goal},\text{ and } v_\text{goal}$ are the goal x-position, y-position, and velocities, ${p_x}_k, {p_y}_k, \text{ and }v_k$ are the x-position, y-position, and velocities of the nominal point $z_k$ while ${p_x}_i, \text{ and }{p_y}_i$ are the x and y positions of points from $\tilde{\mathcal{D}}$, and $G$ is a weight for the goal node (effectively creating G points at the goal). Lastly, $\beta$ is a sharpness parameter to ensure that we minimize the distance to a particular point in $\tilde{\mathcal D}$ as opposed to minimizing the distance to all points. The inclusion of the goal point enabled us to provide a high weight to the cost without hinder the MPC optimizer from reaching the goal. Since our goal is to minimize the distance to the representative points by minimizing the cost function we scale the internal sum by $-1$. For efficient implementation we included this as a part of the quadratic cost, thus requiring the outer-exponential to provide a lower-bound to the cost. In our experiments we choose $L=800$, $G=200$, and $\beta = 3$, and our total cost was $J_\text{total}(\Z,\V) = J(\Z,\V) + J_{\mathcal{X}_\textrm{f}}(\Z,\V) + 1000J_\text{active}(\Z,\V)$. For $Q_f, Q, \text{ and } R$ we used the same parameters as the remaining Dubins' car experiments and let $Q_s = 0.1I_4$. 

\section{Theorem \ref{thm:cov_gap_thm} Analysis}\label{app:theoretical_expansion}

\begin{figure}[htbp]
    \centering
    \includegraphics[width=\textwidth]{figs/weighted_conformal_final_figure.pdf}
    \caption{\textbf{Toy Example.} We plot the effect of the calibration set size, $N_\textrm{calib}$, and compare the weightage provided to points using the exponential weight decay function at different $\rho$ values. The results are averaged across 100 trials of randomly sampling $N_\textrm{calib}$ points in the unit circle, with an example of the calibration data location in the top row. The second row (a histogram) shows that, under uniform sampling on the unit circle, the number of points at each radius (Distance r) increases linearly. The remaining rows provide the weight contribution of each histogram bin by normalizing the weights (as discussed in Section~\ref{sec:cp_err_bound}) and summing the normalized weights for each bin. We observe that smaller $\rho$ values place more emphasis on points closer to the test point, and the weight of the test point, represented by a point mass at $\infty$ in the empirical score distribution, increases. Lastly, in the plot, red denotes the test point, while orange highlights the bin of points closest to it.}
    \label{fig:point_distribution}
\end{figure}

For completeness we restate the \textit{Coverage Gap} bound in Theorem \ref{thm:cov_gap_thm}: 
$$\text{Coverage Gap}\le \underbrace{\sum_{i=1}^{N_\textrm{calib}} \left( \frac{\rho^{d_i}}{1 + \sum_{j=1}^{N_\textrm{calib}} \rho^{d_j}} \right) \cdot 2\epsilon \cdot d_i}_{\text{Tight Bound}} \le \underbrace{2 \epsilon \left[  \frac{d_1 }{1 - \rho^{d_\text{min}}} + \frac{d_\text{max}\rho^{d_\text{min}}}{(1-\rho^{d_\text{min}})^2}\right]}_{\text{Interpretable Bound}}$$
As discussed in Section \ref{sec:theory}, an increase in the density of state-control calibration points reduces $d_\textrm{min}$, thereby increasing the (interpretable) coverage gap bound in \eqref{eq:thm3_bound}. We can reduce the bound by reducing $\rho$. To build intuition, on the results, we consider the following toy example: we have calibration points uniformly sampled from a unit circle with a test point located at the origin, $(0,0)$, such that the distribution of the non-conformity score is spatially varying (i.e. $\tvd \neq 0$). Without defining the distribution, Figure \ref{fig:point_distribution} demonstrates that at smaller $\rho$ values, greater emphasis is provided to non-conformity scores located near the test point, with smaller TV distances to the test point. Hence, the coverage gap decreases for a fixed $N_\textrm{calib}$ as $\rho$ decreases. The interpretable bound also implies that when $d_\textrm{min}$ is reduced, the coverage gap increases, which can be compensated for by requiring $\rho$ to decrease. Figure \ref{fig:point_distribution}, illustrates that the coverage gap increases at larger $d_\textrm{min}$ stemming from increased data density and $N_\textrm{calib}$, due to a higher weightage being provided to the test point, represented as a point-mass at $\infty$ in the empirical weighted non-conformity score distribution, resulting in a larger $1-\alpha$ quantile value. The larger quantile value then increases coverage, thereby reducing the coverage gap. 


To understand the tightness of the derived bound, we use a gamma distribution to randomly sample non-conformity scores such that $s(x_i,y_i) \sim S(x_i,y_i)= \Gamma(\alpha=5 - \frac{1}{3}(\sqrt{x_i^2+y_i^2})^{0.8}, \theta=2.0)$, where $\mathbb{E}[s(x_i,y_i)] = 10 - \frac{2}{3}(\sqrt{x_i^2+y_i^2})^{0.8}$ and $\textrm{Var}[s(x_i,y_i)]=20 - \frac{4}{3}(\sqrt{x_i^2+y_i^2})^{0.8}$. To estimate $\epsilon$, we randomly sampled several pairs of points in the unit circle, then computed $\frac{\tvd (S(x_1,y_1),S(x_2,y_2))}{||(x_1,y_1)-(x_2,y_2)||}$ and took the maximum value computed to be $\epsilon$. Then, we sampled several $N_\textrm{calib}$ values and $\rho$ values on a log scale and computed the coverage across 100 trials. In each trial, we randomly sampled $N_\textrm{calib}$ points on the unit circle, then sampled a non-conformity score from the distribution defined by the coordinates of each sampled point. Since we know the distribution, we can directly compute the TV distance \eqref{eq:init_barber} from \citet{barber2023conformal}. We also computed the tight bound in \eqref{eq:thm3_bound} using the estimated $\epsilon$ value. 

\begin{figure}[htbp]
    \centering
    \includegraphics[width=\textwidth]{figs/weighted_conformal_bounds_extended.pdf}
    \caption{\textbf{Toy Example.} We plot the empirical coverage, along with the theoretical bounds from \eqref{eq:init_barber}, \citet{barber2023conformal}, and our derived bound (labeled) as CP-SLS cov in \eqref{eq:thm3_bound} for $\rho$ values selected on a log scale and several $N_\textrm{calib}$ values. We observe that the lower bounds and empirical coverage increase as $\rho$ decreases, and that the bounds increase more rapidly for smaller $\rho$ values.}
    \label{fig:toy_example_empirical}
\end{figure}

The results in Figure \ref{fig:toy_example_empirical} demonstrate that the lower bound on coverage ($=1-\alpha - \textit{Coverage Gap}$) increases for large $\rho$ as well as the empirically observed coverage. Additionally, the tightness of our derived bound, relative to \citet{barber2023conformal}, improves for smaller $\rho$ values. Similarly, we find the lower-bound increases for fixed $\rho$ values as $N_\textrm{calib}$ increases. While these results suggest a smaller $\rho$ is optimal, for coverage, it is essential to consider the utility sacrifice. Smaller $\rho$ values lead to larger, potentially infinite, values being computed as the $1-\alpha$ quantile, thereby making the uncertainty set harder to use and potentially impractical.  

\newpage
\section{Additional Results}\label{app:additional_results}
% \begin{figure}[h]
%     \centering
%     \includegraphics[width=\textwidth]{figs/ellip_ball_comparison_overlay.pdf}
%     \caption{\textbf{In Domain Car.} We plot the minimum distance of the prediction error to the edge of the ellipsoid. Both the fixed ball and adaptive ellipsoid method empirically satisfy the coverage guarantee and remain inside over $99.33\% = (1- \frac{0.1}{15})\times100\%$ of the time.}
%     \label{fig:ellip_in_domain}
% \end{figure}

% \begin{figure}[h]
%     \centering
%     \includegraphics[width=\textwidth]{figs/run_9_friction_car.pdf}
%     \caption{\textbf{Friction Car.} We provide an example run of the friction area where the adaptive ellipsoid approach is successful in avoiding the obstacle. The remaining approaches fail because they do not maintain sufficient obstacle proximity and thus crash into the obstacle.}
%     \label{fig:second_run_figure}
% \end{figure}
\begin{figure}[!h] % 'p' suggests a dedicated page for floats
    \centering
    
    % First Figure
    \begin{minipage}{1.0\textwidth}
        \centering
    \includegraphics[width=\textwidth]{figs/ellip_ball_comparison_overlay.pdf}
    \caption{\textbf{In Domain Car.} We plot the minimum distance of the prediction error to the edge of the ellipsoid. Both the fixed ball and adaptive ellipsoid method empirically satisfy the coverage guarantee and remain inside over $99.33\% = (1- \frac{0.1}{15})\times100\%$ of the time.}
    \label{fig:ellip_in_domain}
    \end{minipage}

    % \vspace{5mm} % Adjust the vertical gap between the two images
    
    % Second Figure
    \begin{minipage}{1.0\textwidth}
        \centering
    \includegraphics[trim={0cm 0.5cm 0cm 1.3cm}, clip, width=\textwidth]{figs/run_9_friction_car.pdf}
    \vspace{-3cm}
    \caption{\textbf{Friction Car.} We provide an example run of the friction area where the adaptive ellipsoid approach is successful in avoiding the obstacle. The remaining approaches fail because they do not maintain sufficient obstacle proximity and thus crash into the obstacle.}
    \label{fig:second_run_figure}
    \end{minipage}
\end{figure}
\clearpage

\begin{table}[h]
    \small
  \centering
  \caption{\textbf{Active Uncertainty.} The computation times (ms), and prediction errors when planning with vanilla-MPC, the adaptive ellipsoid with and without active uncertainty. }
  \label{tab:active_uncert}
  \begin{tabular}{l|ccc}
    \hline
    \textbf{Model} & \textbf{V-MPC} & \textbf{CP-Ellipsoid} & \textbf{CP-Ellipsoid + Active Uncertainty} \\
    \hline
    \textbf{Time (ms)} & 56.1 $\pm$ 1.1 & 132.0 $\pm$ 9.9 & 131.4 $\pm$ 8.2 \\
    \hline
    \textbf{Pred. Err. ($L_2$) ($\times 10^{-2}$)} & 2.82 $\pm$ 7.85 & 2.83 $\pm$ 7.89 & .34 $\pm$ .59 \\
    \hline
  \end{tabular}
\end{table}

% \begin{figure}[h]
%     \centering
%     \includegraphics[width=\textwidth]{figs/run_9_c_shape_car.pdf}
%     \caption{\textbf{Active Uncertainty.} We provide an example run of the rollouts when using V-MPC and CP-Ellipsoid with active uncertainty. We find that the active uncertainty mostly avoid the OOD region, and thus has a smoother more certain trajectory.}
%     \label{fig:rollout_example_1_active_uncert}
% \end{figure}

% \begin{figure}[h]
%     \centering
%     \includegraphics[width=\textwidth]{figs/run_2_c_shape_car.pdf}
%     \caption{\textbf{Active Uncertainty.} We provide another example run of the rollouts when using V-MPC and CP-Ellipsoid with active uncertainty. Similar to Figure \ref{fig:rollout_example_1_active_uncert} the active uncertainty rollout is smoother by mitigating time in the OOD region. Although, the active uncertainty approach has more steps in the OOD region than Figure \ref{fig:rollout_example_1_active_uncert}, it is near the in distribution region and thus encounter limited OOD errors in the rollout.}
%     \label{fig:rollout_example_2_active_uncert}
% \end{figure}
\begin{figure}[p] % 'p' suggests a dedicated page for floats
    \centering
%     \begin{minipage}{1.0\textwidth}
%     \small
%   \centering
%   \label{tab:active_uncert}
%     \captionof{table}{\textbf{Active Uncertainty.} The computation times (ms), and prediction errors when planning with vanilla-MPC, the adaptive ellipsoid with and without active uncertainty. }
% \vspace{-2mm}
%   \begin{tabular}{l|ccc}
%     \hline
%     \textbf{Model} & \textbf{V-MPC} & \textbf{CP-Ellipsoid} & \textbf{CP-Ellipsoid + Active Uncertainty} \\
%     \hline
%     \textbf{Time (ms)} & 56.1 $\pm$ 1.1 & 132.0 $\pm$ 9.9 & 131.4 $\pm$ 8.2 \\
%     \hline
%     \textbf{Pred. Err. ($L_2$) ($\times 10^{-2}$)} & 2.82 $\pm$ 7.85 & 2.83 $\pm$ 7.89 & .34 $\pm$ .59 \\
%     \hline
%   \end{tabular}
%     \end{minipage}
    
    \begin{minipage}{1.0\textwidth}
        \centering
    \includegraphics[trim={0cm 2cm 0cm 2cm}, clip,width=0.95\textwidth]{figs/run_9_c_shape_car.pdf}
    \vspace{-1.8cm}
    \caption{\textbf{Active Uncertainty.} We provide an example run of the rollouts when using V-MPC and CP-Ellipsoid with active uncertainty. We find that the active uncertainty approach mostly avoids the OOD region, and thus has a smoother, more certain trajectory.}
    \label{fig:rollout_example_1_active_uncert}
    \end{minipage}

    % Second Figure
    \begin{minipage}{1.0\textwidth}
        \centering
    \includegraphics[trim={0cm 2cm 0cm 2cm}, clip,width=0.95\textwidth]{figs/run_2_c_shape_car.pdf}
    \vspace{-1.8cm}
    \caption{\textbf{Active Uncertainty.} We provide another example run of the rollouts when using V-MPC and CP-Ellipsoid with active uncertainty. Similar to Figure \ref{fig:rollout_example_1_active_uncert}, the active uncertainty rollout is smoother by mitigating time in the OOD region. Although the active uncertainty approach has more steps in the OOD region than shown in Figure \ref{fig:rollout_example_1_active_uncert}, it is near the in-distribution region and thus encounters limited OOD errors during the rollout.}
    \label{fig:rollout_example_2_active_uncert}
    \end{minipage}
\end{figure}
\begin{figure}[!h]
    \centering
    \includegraphics[width=\textwidth]{figs/quad_trajectory_light_5.pdf}
    \caption{\textbf{Quadcopter.} We provide another example run of the quadcopter rollout, with a different initial condition than in Figure \ref{fig:quad_trajectory}. Again, we observe the quadcopter is successful in avoiding the obstacle when using CP-Ellipsoid (unlike V-MPC) when entering the OOD region. }
    \label{fig:quad_traj_additional_result}
\end{figure}
\end{document}




